\documentclass{article}
\usepackage{anyfontsize}
\usepackage[UTF8]{ctex} % 如果需要支持中文
\usepackage{fontspec} 
% \setCJKmainfont{SourceHanSansCN-Light}[
%     BoldFont=SourceHanSansCN-Medium
% ]

\usepackage[a4paper, left=0.1in, right=0.1in, top=0.05in, bottom=0.05in]{geometry} % 设置四边页边距为0.1英寸{geometry} % 设置页边距为0.5英寸
\usepackage{enumitem} % 用于自定义列表
\usepackage{amsmath}  % 数学公式支持
\usepackage{tikz}
\usetikzlibrary{arrows.meta, positioning}
\setlength{\parindent}{0pt} % 不缩进
\usepackage{etoolbox} % 用于列表操作
%调整类代码
\usepackage{comment}
\usepackage{tocloft} % 用于定制目录 样式
\usepackage{hyperref} %不需要目录盒子注释这
\usepackage{ifthen}
\usepackage{tikz}
\usetikzlibrary{arrows.meta}
\usepackage{titletoc}
\usepackage{graphicx}
\usepackage{float}

\usepackage{amssymb}  % 提供数学符号，包括\therefore
%目录设定
%快捷键 CMD + / 注释
%  \renewcommand{\cftdot}{} % 隐藏目录中的页码
%  \cftpagenumbersoff{section}
%  \cftpagenumbersoff{subsection}
%  \makeatletter
%  \renewcommand{\@pnumwidth}{0em}  % 设置页码宽度为0
%  \renewcommand{\@tocrmarg}{0em}   % 设置右边距为0
%  \makeatother
 



\usepackage{environ}

\newif\ifshowcontent  % 定义一个条件判断变量
\showcontenttrue    % 设置为 false，隐藏所有 question 环境的内容

\newcounter{question} % 题目计数器
\setcounter{question}{0}
\NewEnviron{question}{%
    \ifshowcontent
        \par\stepcounter{question}\textbf{\thequestion.} \BODY\par % 如果显示内容，则输出
    \fi
}

\NewEnviron{choices}{%
    \ifshowcontent
        \begin{enumerate}[label=\Alph*., leftmargin=*]
            \BODY
        \end{enumerate}\par
    \fi
}

\NewEnviron{correctanswer}{%
    \ifshowcontent
        \textbf{答案:}\BODY\par
    \fi
}



\newenvironment{explanation}
{\textbf{解析:}\par}
{\par}



% 新增考察方面命令
\newcommand{\checklist}[1]{
    \textbf{考察:} #1 \par % 在题目下方显示考察方面
    \addcontentsline{toc}{subsection}{题号\thequestion 考察: #1} % 将考察内容添加到目录
    \vspace{2em} % 添加1em的垂直空间
}

\graphicspath{{images/}} %统一


\begin{document}
 %\fontsize{22}{31}\selectfont
 \tableofcontents % 添加目录
\newpage
%\pagestyle{empty}






\section{考点1:衍生品风险}
\setcounter{question}{0}

\begin{question}
    Which of the following statements best describes the benefit of netting risk exposures? The benefits of netting are realized when:
    \\ 哪一项陈述最准确地描述了净额结算风险敞口的好处？净额结算的好处是在以下情况下实现的：
\end{question}

\begin{choices}
    \item Marked-to-market (MtM) values have high structural correlations for two trades\\
    两笔交易的Marked-to-market (MtM)值具有高结构相关性
    \item Marked-to-market (MtM) values have opposite signs for two trades\\
    两笔交易的Marked-to-market (MtM)值具有相反的符号
    \item Expected exposure (EE) values are minimal\\
    预期敞口（EE）值最小
    \item Expected future exposure (EFE) values have zero correlation\\
    预期未来敞口（EFE）值具有零相关性
\end{choices}

\begin{correctanswer}
    B
\end{correctanswer}

\begin{explanation}
    \textbf{A选项错误。}
    \begin{itemize}
        \item 高结构相关性并不意味着可以实现净额结算的好处。即使两笔交易高度相关，如果它们的MtM值符号相同，净额结算也无法减少风险敞口。
    \end{itemize}

    \textbf{B选项正确。}
    \begin{itemize}
        \item 当两笔交易的MtM值具有相反的符号时，净额结算可以显著减少风险敞口。例如，如果一笔交易的价值为+100，另一笔为-80，净额结算后的敞口仅为20，而不是180。
    \end{itemize}

    \textbf{C选项错误。}
    \begin{itemize}
        \item 预期敞口（EE）值最小化并不是净额结算的必要条件。即使EE值较大，只要MtM值符号相反，净额结算仍然可以带来好处。
    \end{itemize}

    \textbf{D选项错误。}
    \begin{itemize}
        \item 预期未来敞口（EFE）值的相关性为零并不意味着可以实现净额结算的好处。净额结算主要依赖于MtM值的符号，而不是EFE值的相关性。
    \end{itemize}
\end{explanation}

\checklist{净额结算（netting）的实现条件}


\begin{question}
    XYZ Hedge Fund wants to get exposure to a high-yield pool of commercial loans without actually investing in the loans. It wants a leverage ratio of 8.0. If the hedge fund is willing to invest \$40 million in this investment, which credit derivative is best for them and what is their expected return given that the reference asset earns LIBOR plus 300 basis points, the counterparty earns LIBOR plus 160 basis points, and the required collateral earns 4.0\%?
    \\ XYZ对冲基金想要获得对高收益商业贷款池的敞口，而不实际投资于这些贷款。它希望杠杆率为8.0。如果对冲基金愿意投资\$40百万，哪一种信用衍生品最适合他们，他们的预期回报是多少，假设参考资产赚取LIBOR加上300个基点，交易对手赚取LIBOR加上160个基点，所需的抵押品赚取4.0\%？
\end{question}

\begin{choices}
    \item Total return swap with a 14.80\% return\\
    总收益互换，回报率为14.80\%
    \item Asset-backed credit-linked note with a 12.40\% return\\
    资产支持信用关联票据，回报率为12.40\%
    \item Total return swap with a 12.40\% return\\
    总收益互换，回报率为12.40\%
    \item Asset-backed credit-linked note with a 14.80\% return\\
    资产支持信用关联票据，回报率为14.80\%
\end{choices}

\begin{correctanswer}
    D
\end{correctanswer}
\begin{explanation}
\textbf{A选项错误。}
\begin{itemize}
    \item 总收益互换（Total Return Swap）的回报率计算方式不同，通常回报率低于资产支持信用关联票据。在总收益互换中，对冲基金作为总收益接收方，回报率通常为12.40\%，而不是14.80\%。
\end{itemize}

\textbf{B选项错误。}
\begin{itemize}
    \item 资产支持信用关联票据的回报率计算。具体计算看D选项解析。
\end{itemize}

\textbf{C选项错误。}
\begin{itemize}
    \item 总收益互换的回报率为12.40\%，低于资产支持信用关联票据的14.80\%。此外，总收益互换不是实现杠杆化信用敞口的最佳工具，因为资产支持信用关联票据可以通过抵押品和利差获得更高的回报。
\end{itemize}

\textbf{D选项正确。}
\begin{itemize}
    \item 对冲基金应该选择资产支持信用关联票据，因为：
    \begin{itemize}
        \item 资产支持信用关联票据允许对冲基金通过抵押品投资获得杠杆化的信用敞口，而不需要直接投资于商业贷款。
        \item 杠杆率8.0意味着名义价值为$\$320\text{ million}$（$8.0 \times \$40\text{ million}$）。
        \item 对冲基金从\$40 million的抵押品中获得4.0\%的收益，即\$1.6 million。
        \item 对冲基金从\$320 million的名义价值中获得140个基点的利差（300 bps - 160 bps），即\$4.48 million。
        \item 总回报率：\[\frac{\$1.6\text{ million} + \$4.48\text{ million}}{\$40\text{ million}} = 14.80\%\]
    \end{itemize}
\end{itemize}
\end{explanation}


\checklist{资产支持信用关联票据的回报计算}


\begin{question}
    A portfolio manager enters into a total return swap as the total return receiver. Under which of the following situations, would the portfolio manager be required to make a net outlay to the counterparty?
    \\ 一个投资组合经理作为总收益接收方进入总收益互换。在以下情况下，投资组合经理需要向交易对手支付净额？
\end{question}

\begin{choices}
    \item If the transaction was initiated as a hedge. Then no outlay was required\\
    如果交易作为对冲工具发起，则不需要支付净额
    \item If there were a capital gain on the reference asset\\
    如果参考资产出现资本利得，则需要支付净额
    \item If the market value of the reference asset decreased significantly\\
    如果参考资产的市场价值显著下降，则需要支付净额
    \item If the spread between the reference asset and the benchmark asset changed\\
    如果参考资产与基准资产之间的利差发生变化，则需要支付净额
\end{choices}

\begin{correctanswer}
    C
\end{correctanswer}

\begin{explanation}
    \textbf{A选项错误。}
    \begin{itemize}
    \item 交易是否作为对冲工具发起与是否需要支付净额无关。总收益互换的支付义务取决于参考资产的表现，而不是交易的目的。
    \end{itemize}

    \textbf{B选项错误。}
    \begin{itemize}
    \item 如果参考资产出现资本利得，总收益接收方（投资组合经理）将收到支付，而不是需要支付。总收益接收方在资产升值时获得收益，在资产贬值时承担损失。
    \end{itemize}

    \textbf{C选项正确。}
    \begin{itemize}
    \item 当参考资产的市场价值显著下降时，总收益接收方需要向对手方支付净差额。这是因为总收益接收方同时面临信用风险和市场风险，必须承担参考资产价值下降带来的损失。
    \end{itemize}

    \textbf{D选项错误。}
    \begin{itemize}
    \item 参考资产与基准资产之间的利差变化本身并不会导致总收益接收方需要支付净额。支付义务主要取决于参考资产的总收益表现。
    \end{itemize}
\end{explanation}

\checklist{总收益互换（Total Return Swap）的支付义务}


\begin{question}
    Which of the following statements is not correct regarding total return swaps (TRS)?
    \\ 哪一项陈述关于总收益互换（TRS）是不正确的？
\end{question}

\begin{choices}
    \item A TRS is designed to mirror the return on an underlying asset like a loan, stock, or even a portfolio of assets\\
    总收益互换（TRS）设计目的是复制标的资产（如贷款、股票或资产组合）的收益
    \item The payer pays any depreciation in the underlying asset to the receiver\\
    支付方支付标的资产的贬值部分给接收方
    \item The payer pays any dividends or interest received to the receiver\\
    支付方支付收到的任何股息或利息给接收方
    \item The receiver is creating a synthetic long position in the underlying\\
    接收方创建一个标的资产的合成多头头寸
\end{choices}

\begin{correctanswer}
    B
\end{correctanswer}

\begin{explanation}
    \textbf{A选项说法恰当。}
    \begin{itemize}
        \item 总收益互换（TRS）的设计目的就是复制标的资产（如贷款、股票或资产组合）的收益。
    \end{itemize}

    \textbf{B选项说法不恰当。}
    \begin{itemize}
        \item 在总收益互换中，支付方（payer）只负责支付标的资产的升值部分以及相关的股息或利息。而接收方（receiver）需要向支付方支付标的资产的贬值部分。总收益互换同时转移了信用风险和市场风险。
    \end{itemize}

    \textbf{C选项说法恰当。}
    \begin{itemize}
        \item 支付方确实需要将收到的任何股息或利息支付给接收方，这是总收益互换的基本机制之一。
    \end{itemize}

    \textbf{D选项说法恰当。}
    \begin{itemize}
        \item 接收方通过总收益互换实际上创建了一个对标的资产的合成多头头寸，因为接收方获得了标的资产的所有收益（包括升值和收入），同时也承担了所有损失。
    \end{itemize}
\end{explanation}

\checklist{总收益互换（Total Return Swap）的支付义务}


\begin{question}
    In a total return swap, the:
    \\ 在总收益互换中，哪一项是正确的？
\end{question}

\begin{choices}
    \item Total return payer bears the risk of holding the reference asset\\
    总收益支付方承担持有参考资产的风险
    \item Total return payer transfers the risk of holding the reference asset to the total return receiver\\
    总收益支付方将持有参考资产的风险转移给总收益接收方
    \item Total return receiver transfers the risk of holding the reference asset to the total return payer\\
    总收益接收方将持有参考资产的风险转移给总收益支付方
    \item Credit risk of the buyer and seller cancel each other out
    \item 买卖双方的信用风险相互抵消
\end{choices}

\begin{correctanswer}
    B
\end{correctanswer}

\begin{explanation}
    \textbf{A选项错误。}
    \begin{itemize}
    \item 总收益支付方（payer）并不承担持有参考资产的风险。相反，他们通过总收益互换将风险转移给了总收益接收方。
    \end{itemize}
    \textbf{B选项正确。}
    \begin{itemize}
    \item 在总收益互换中，总收益支付方将参考资产的资本利得和损失转移给了总收益接收方。总收益互换是一种信用衍生产品，它将标的资产（如贷款、债券或资产组合）的总收益与LIBOR加上信用价差进行交换。总收益包括利息收入和资产价值变动。总收益互换不仅可以管理信用风险，还可以管理市场风险。
    \end{itemize}
    \textbf{C选项错误。}
    \begin{itemize}
    \item 总收益接收方（receiver）并不将风险转移给支付方，而是从支付方那里接收风险。接收方通过支付LIBOR加上价差来获得参考资产的总收益。
    \end{itemize}
    \textbf{D选项错误。}
    \begin{itemize}
    \item 买卖双方的信用风险并不会相互抵消。总收益互换中的信用风险主要存在于交易对手方之间，而不是相互抵消。
    \end{itemize}
\end{explanation}

\checklist{总收益互换（Total Return Swap）的支付义务}

\begin{question}
    Risk Averse Bank (RAB) has made a loan of USD 120 million at 7.5\% per annum. RAB wants to enter into a total return swap under which it will pay the interest on the loan plus the change in the mark-to-market value of the loan and in exchange, RAB will get LIBOR + 40bps. Settlement payments are made annually. What is the cash flow for RAB on the first settlement data if the mark-to-market value of the loan falls by 2.5\% and LIBOR is 5.5\%?
    \\ 风险厌恶银行（RAB）以7.5\%的年利率发放了USD 120百万的贷款。RAB希望进入总收益互换，在其中它将支付贷款利息加上贷款市值的变化，作为回报，RAB将获得LIBOR加上40个基点。结算支付每年进行一次。如果贷款市值下降了2.5\%，LIBOR为5.5\%，RAB在第一次结算时的现金流是多少？
\end{question}

\begin{choices}
    \item Net inflow of USD 0.4 million\\
    净流入USD 0.4百万
    \item Net outflow of USD 0.4 million\\
    净流出USD 0.4百万
    \item Net inflow of USD 2.0 million\\
    净流入USD 2.0百万
    \item Net outflow of USD 2.0 million\\
    净流出USD 2.0百万
\end{choices}

\begin{correctanswer}
    A
\end{correctanswer}

\begin{explanation}


    \textbf{步骤1：理解总收益互换的现金流结构}
    \begin{itemize}
        \item RAB作为支付方（payer），需要支付：贷款利息 + 贷款市值变化
        \item RAB作为接收方（receiver），收到：LIBOR + 40 bps
    \end{itemize}

    \textbf{步骤2：计算RAB需要支付的金额}
    \begin{itemize}
        \item 贷款利息：$\$120\text{ million} \times 7.5\% = \$9\text{ million}$
        \item 贷款市值下降（贬值）：$\$120\text{ million} \times 2.5\% = \$3\text{ million}$
        \item RAB需要支付的总收益：$\$9\text{ million} - \$3\text{ million} = \$6\text{ million}$
        \item 注意：市值下降意味着RAB需要向对手方支付贬值部分
    \end{itemize}

    \textbf{步骤3：计算RAB收到的金额}
    \begin{itemize}
        \item RAB收到的浮动利率：$\text{LIBOR} + 40\text{ bps} = 5.5\% + 0.4\% = 5.9\%$
        \item RAB收到的金额：$\$120\text{ million} \times 5.9\% = \$7.08\text{ million}$
    \end{itemize}

    \textbf{步骤4：计算净现金流}
    \begin{itemize}
        \item 净流入：$\$7.08\text{ million} - \$6\text{ million} = \$1.08\text{ million}$
    \end{itemize}
\end{explanation}

\checklist{总收益互换（Total Return Swap）的现金流计算}


\begin{question}
    The credit exposure of an interest rate swap differs from that of a bond in that:

    I. The swap can be terminated by novation

    II. The principal amount of the swap is not at risk

    III. Swaps benefit from higher recovery rates

    IV. The full coupon amounts in the swap are not at risk
    \\ 利率互换的信用风险与债券的信用风险不同，因为：

    I. 互换可以通过更新终止

    II. 互换的名义本金不面临信用风险

    III. 互换受益于更高的回收率

    IV. 互换中的全部息票金额不面临风险，只有净差额面临风险
\end{question}

\begin{choices}
    \item I and III
    \item II and IV
    \item II, III and IV
    \item I, II, III and IV
\end{choices}

\begin{correctanswer}
    B
\end{correctanswer}

\begin{explanation}
    \textbf{陈述I不正确。}
    \begin{itemize}
        \item 虽然互换可以通过更新终止（novation），但这并不是互换信用风险区别于债券的主要特征。债券也可以通过提前偿还等方式终止，因此这不能作为互换与债券信用风险的主要区别。
    \end{itemize}

    \textbf{陈述II正确。}
    \begin{itemize}
        \item 在利率互换中，名义本金不需要交换，双方只交换利息差额。因此，本金金额不面临信用风险，这与债券不同，债券的本金在到期时面临信用风险。
    \end{itemize}

    \textbf{陈述III不正确。}
    \begin{itemize}
        \item 互换的回收率并不一定高于债券。回收率主要取决于交易的优先级、是否有抵押品、债务人的信用状况等因素，而不是取决于工具类型本身。因此，互换并不一定享有更高的回收率。
    \end{itemize}

    \textbf{陈述IV正确。}
    \begin{itemize}
        \item 在利率互换中，双方只需要交换固定利率和浮动利率之间的净差额，而不是全部的息票金额。例如，如果固定利率为5\%，浮动利率为3\%，则只需要支付2\%的净差额，而不是支付5\%的固定利率。因此，互换中的全部息票金额不面临风险，只有净差额面临风险。
    \end{itemize}
\end{explanation}

\checklist{利率互换（Interest Rate Swap）与债券的信用风险比较}


\begin{question}
    An investor can gain the exposure and return of an underlying loan by:
    \\ 投资者可以通过以下方式获得标的贷款的风险敞口和收益：
\end{question}

\begin{choices}
    \item Being the credit risk buyer in a total return swap\\
    作为总收益互换中的信用风险买方
    \item Being the credit risk seller in a total return swap\\
    作为总收益互换中的信用风险卖方
    \item Entering into an interest-rate swap with a BB credit\\
    进入一个与BB信用等级的利率互换
    \item Pledging compensating balances with the lender bank\\
    向贷款银行提供补偿性余额
\end{choices}

\begin{correctanswer}
    A
\end{correctanswer}

\begin{explanation}
    \textbf{A选项正确。}
    \begin{itemize}
    \item 在总收益互换中，信用风险卖方（总收益支付方）将标的资产的总收益支付给信用风险买方（总收益接收方）。通过成为总收益互换中的信用风险买方，投资者可以获得标的贷款的风险敞口和收益，而不需要实际持有该贷款。
    \end{itemize}

    \textbf{B选项错误。}
    \begin{itemize}
    \item 作为总收益互换中的信用风险卖方，投资者实际上是转移了标的贷款的风险和收益，而不是获得这些风险和收益。信用风险卖方需要支付标的资产的总收益，并收取LIBOR加上价差。
    \end{itemize}

    \textbf{C选项错误。}
    \begin{itemize}
    \item 与BB级信用的对手方进行利率互换并不能获得标的贷款的风险敞口和收益。利率互换只涉及利率的交换，不涉及信用风险的转移。
    \end{itemize}

    \textbf{D选项错误。}
    \begin{itemize}
    \item 向贷款银行提供补偿性余额并不能获得标的贷款的风险敞口和收益。补偿性余额只是作为贷款的担保，并不涉及风险和收益的转移。
    \end{itemize}
\end{explanation}

\checklist{总收益互换（Total Return Swap）的信用风险买方可以获得标的贷款的风险敞口和收益}


\begin{question}
    A standard synthetic CDO (basket credit default swap) references a portfolio of twelve (12) individual corporate names. Assume the following:
    \begin{itemize}
    \item The total reference notional (basket notional) is X, and the term is Y years
    \item The reference notional per individual reference credit name is X/12 (i.e. equal weight per name)
    \item The default correlations between the individual reference credit names are all equal to one (1.0)
    \item The single-name credit default swap (CDS) spread for each individual reference credit name is 120 basis points, for a term of Y years
    \item The assumed recovery rate on default for all individual reference credits is zero in all cases
    \item The synthetic CDO comprises two tranches, a 40\% junior tranche priced at a spread J, and a 60\% senior tranche priced at spread S
    \end{itemize}
    All else held constant, if the default correlations between the individual reference credit names are reduced from 1.0 to 0.6, what is the effect on the relationship between the junior tranche spread J and the senior tranche spread S?
    \\ 一个标准合成CDO（篮子信用违约互换）参考了一个包含十二个（12）个独立企业名称的组合。假设：
    \begin{itemize}
    \item 参考名义总额（篮子名义总额）为X，期限为Y年
    \item 每个独立参考信用名称的参考名义为X/12（即每个名称权重相等）
    \item 各独立参考信用名称之间的违约相关性均为1.0
    \item 每个独立参考信用名称的单名信用违约互换（CDS）利差为120个基点，期限为Y年
    \item 所有独立参考信用名称的违约回收率为零
    \item 合成CDO包含两个层级，一个40\%的次级层级，定价利差为J，一个60\%的高级层级，定价利差为S
    \end{itemize}
    所有其他条件不变，如果各独立参考信用名称之间的违约相关性从1.0降低到0.6，次级层级利差J和高级层级利差S之间的关系会发生什么变化？
\end{question}

\begin{choices}
    \item The relationship remains the same\\
    关系保持不变
    \item S increases relative to J\\
    S相对于J增加
    \item J increases relative to S\\
    J相对于S增加
    \item The effect cannot be determined given the data supplied\\
    由于提供的数据，无法确定影响
\end{choices}

\begin{correctanswer}
    C
\end{correctanswer}

\begin{explanation}
    \textbf{A选项错误。}
    \begin{itemize}
    \item 当违约相关性从1.0降低到0.6时，层级之间的关系会发生变化。相关性变化会影响不同层级的风险分配。
    \end{itemize}

    \textbf{B选项错误。}
    \begin{itemize}
    \item 当违约相关性降低时，高级层（senior tranche）的利差不会相对于次级层（junior tranche）增加。相反，次级层的利差会增加。
    \end{itemize}

    \textbf{C选项正确。}
    \begin{itemize}
    \item 当违约相关性从1.0降低到0.6时：
    \begin{itemize}
        \item 在相关性为1.0时，只有两种可能的结果：要么所有发行人同时违约，要么没有发行人违约。
        \item 当相关性降低到0.6时，会出现部分发行人违约的情况。
        \item 这些部分违约的情况对次级层的影响大于对高级层的影响。
        \item 虽然整体平均违约率保持不变，但次级层承担了更多的风险。
        \item 因此，次级层的利差相对于高级层会增加，因为整个组合的总利差必须保持在120个基点。
    \end{itemize}
\end{itemize}
    \textbf{D选项错误。}
    \begin{itemize}
    \item 根据提供的数据，可以确定违约相关性降低对层级利差关系的影响。次级层的利差会相对于高级层增加。
    \end{itemize}
\end{explanation}

\checklist{合成CDO（Synthetic CDO）的层级利差关系}


\begin{question}
    Blackstone Credit, Inc., made a loan to a small start-up firm. The firm grew rapidly, and it appeared that Blackstone had made a good credit decision. However, the firm grew too fast and could not sustain the growth. It eventually failed. Blackstone had initially estimated its exposure at default to be \$1,500,000. Because of the firm's rapid growth and resulting increases in the line of credit, Blackstone ultimately lost \$1,850,000. In terms of credit risk, this is an example of:
    \\ Blackstone Credit, Inc.向一家小型初创企业发放了一笔贷款。该企业迅速增长，看起来Blackstone做出了一项良好的信用决策。然而，该企业增长过快，无法持续增长，最终倒闭。Blackstone最初估计的违约风险敞口为\$1,500,000。由于该企业快速增长和由此导致的信贷额度增加，Blackstone最终损失了\$1,850,000。从信用风险的角度来看，这是一个例子：
\end{question}

\begin{choices}
    \item Default on payment for goods or services already rendered\\
    对已提供商品或服务的付款违约
    \item A more severe loss than expected due to a ratings downgrade by a rating agency\\
    由于评级机构降级导致的预期损失更严重
    \item A more severe loss than expected due to a greater than expected exposure at the time of a default\\
    由于违约时的风险敞口大于预期导致的预期损失更严重
    \item A more severe loss than expected due to a lower than expected recovery at the time of a default\\
    由于违约时的回收率低于预期导致的预期损失更严重
\end{choices}

\begin{correctanswer}
    C
\end{correctanswer}

\begin{explanation}
    \textbf{A选项错误。}
    \begin{itemize}
    \item 这个案例涉及的是贷款违约，而不是已提供商品或服务的付款违约。Blackstone提供的是信贷资金，而不是商品或服务。
    \end{itemize}
    \textbf{B选项错误。}
    \begin{itemize}
    \item 这个案例中没有提到评级机构降级的情况。作为一家小型初创企业，该公司可能根本没有信用评级。
    \end{itemize}

    \textbf{C选项正确。}
    \begin{itemize}
    \item 这个案例展示了由于违约时的风险敞口大于最初估计而导致的损失超出预期的情况：Blackstone最初估计的违约风险敞口为150万美元，由于公司快速增长，信贷额度增加，最终损失达到185万美元，比预期多出35万美元。这种损失是由于违约时的实际风险敞口大于最初估计造成的。
    \end{itemize}
    \textbf{D选项错误。}
    \begin{itemize}
    \item 这个案例中没有提到回收率的问题。损失超出预期是由于风险敞口增加，而不是回收率低于预期。
    \end{itemize}
\end{explanation}

\checklist{违约时的风险敞口大于预期导致损失超出预期}


\begin{question}
    ABC Company, domiciled in a country with a strong legal system, is applying for a USD 120,000 loan with an annual interest rate of 6\% to be used exclusively on expanding its business and not to repay current debt. The loan will be secured by the company's factory buildings, which have an appraised value of USD 250,000. The company has an annual operating profit of USD 4,000. Which of the following statements relating to the credit risk of this company can most likely be inferred from the given data?
    \\ ABC公司，位于一个法律体系强大的国家，申请了一笔USD 120,000的贷款，年利率为6\%，用于扩大业务，而不是偿还当前债务。贷款将由公司的工厂建筑提供担保，评估价值为USD 250,000。公司年营业利润为USD 4,000。根据给定数据，哪一项关于该公司信用风险的陈述最有可能被推断出来？
\end{question}

\begin{choices}
    \item The company's capacity to pay is low\\
    公司偿付能力很低
    \item The company's willingness to pay is low\\
    公司还款意愿很低
    \item The loan exhibits a high loss given default\\
    贷款在违约时损失很高
    \item The loan's exposure at default will increase\\
    贷款在违约时风险敞口会增加
\end{choices}

\begin{correctanswer}
    A
\end{correctanswer}

\begin{explanation}
    \textbf{A选项正确。}
    \begin{itemize}
    \item 从给定数据可以推断出公司的偿付能力很低：
    \begin{itemize}
        \item 贷款金额为12万美元，年利率为6\%
        \item 年利息支出为7,200美元（12万 × 6\%）
        \item 公司年营业利润仅为4,000美元
        \item 营业利润不足以支付利息支出，更不用说偿还本金
        \item 这表明公司的偿付能力（capacity to pay）很低
    \end{itemize}
    \end{itemize}
    \textbf{B选项错误。}
    \begin{itemize}
    \item 从给定数据无法推断出公司的还款意愿（willingness to pay）很低。公司申请贷款用于业务扩张而不是偿还现有债务，这表明公司有良好的还款意愿。
    \end{itemize}
    \textbf{C选项错误。}
    \begin{itemize}
    \item 从给定数据无法推断出贷款在违约时会有高损失。抵押物（工厂建筑）的评估价值为25万美元，远高于贷款金额，这表明违约时的损失可能较低。
    \end{itemize}

    \textbf{D选项错误。}
    \begin{itemize}
    \item 从给定数据无法推断出违约时的风险敞口会增加。贷款金额是固定的12万美元，没有信息表明这个金额会随时间增加。
    \end{itemize}
\end{explanation}

\checklist{公司的偿付能力分析}


\begin{question}
    An endowment fund has sold default protection on the most senior tranche of a CDO. If the default correlation between assets held in the CDO decreases sharply from the correlation used in pricing the CDO tranches, assuming everything else is unchanged, how will the position of the endowment fund be impacted?
    \\ 一个捐赠基金出售了CDO最高层级的违约保护。如果CDO中资产的违约相关性从定价CDO层级时使用的相关性急剧下降，假设其他所有条件不变，捐赠基金的持仓将受到什么影响？
\end{question}

\begin{choices}
    \item It will either increase or decrease, depending on the pricing model used and the market conditions\\
    它将增加或减少，取决于使用的定价模型和市场条件
    \item It will gain significant value, since the probability of exercising the protection falls\\
    它将获得显著价值，因为行使保护的概率下降
    \item It will lose significant value, since the protection will gain value\\
    它将失去显著价值，因为保护将获得价值
    \item It will neither gain nor lose value, since only expected default losses matter and correlation does not affect expected default losses\\
    它将既不增加也不减少价值，因为只有预期违约损失事项，相关性不影响预期违约损失
\end{choices}

\begin{correctanswer}
    B
\end{correctanswer}

\begin{explanation}
    \textbf{A选项错误。}
    \begin{itemize}
    \item 违约相关性降低对最高层级（senior tranche）的影响是确定的，不会因为定价模型或市场条件而改变方向。相关性降低总是会增加最高层级的价值。
    \end{itemize}

    \textbf{B选项正确。}
    \begin{itemize}
    \item 当违约相关性降低时，最高层级的价值会增加，原因如下：
    \begin{itemize}
        \item 高相关性意味着如果一个资产违约，CDO中的大量其他资产也会违约
        \item 低相关性意味着如果一个资产违约，对其他资产的违约概率影响很小
        \item 随着相关性降低，发生足够多违约以超过最高层级信用增强的概率也会降低
        \item 因此，出售最高层级保护的投资者（捐赠基金）将获得收益
    \end{itemize}
    \end{itemize}
    \textbf{C选项错误。}
    \begin{itemize}
    \item 这个说法与实际情况相反。违约相关性降低会增加最高层级的价值，而不是降低其价值。
    \end{itemize}

    \textbf{D选项错误。}
    \begin{itemize}
    \item 违约相关性确实会影响CDO的价值，特别是对最高层级的影响。相关性变化会影响违约的集中度，从而影响不同层级的风险。
    \end{itemize}
\end{explanation}

\checklist{合成CDO（Synthetic CDO）的层级利差关系}


\begin{question}
    When considering counterparty credit risk, which of the following financial products has the largest outstanding notional amount in the marketplace?
    \\ 在考虑对手方信用风险时，哪一种金融产品在市场上具有最大的未偿名义金额？
\end{question}

\begin{choices}
    \item Credit default swaps\\
    信用违约互换
    \item Foreign exchange forwards\\
    外汇远期合约
    \item Interest rate swaps\\
    利率互换
    \item Repos and reverse repos\\
    回购和逆回购
\end{choices}

\begin{correctanswer}
    C
\end{correctanswer}

\begin{explanation}
    \textbf{A选项错误。}
    \begin{itemize}
    \item 信用违约互换（CDS）虽然是一个重要的场外衍生品市场，但其名义金额远小于利率互换市场。CDS主要用于信用风险转移，市场规模相对较小。
    \end{itemize}

    \textbf{B选项错误。}
    \begin{itemize}
    \item 外汇远期合约是重要的场外衍生品，但其名义金额也远小于利率互换市场。外汇市场虽然交易活跃，但单个合约的规模通常较小。
    \end{itemize}

    \textbf{C选项正确。}
    \begin{itemize}
    \item 利率互换是场外衍生品市场中规模最大的产品：
    \begin{itemize}
        \item 利率互换占据了场外衍生品市场的大部分份额
        \item 利率互换被广泛用于利率风险管理和资产负债管理
        \item 由于利率风险管理的普遍需求，利率互换的名义金额远大于其他金融产品
    \end{itemize}
    \end{itemize}
    \textbf{D选项错误。}
    \begin{itemize}
    \item 回购和逆回购虽然也是重要的金融产品，但它们属于证券融资交易，而不是场外衍生品。它们的市场规模小于利率互换市场。
    \end{itemize}
\end{explanation}

\checklist{利率互换（Interest Rate Swap）}


\begin{question}
    All of the following could help to reduce the credit exposure on a set of derivative transactions except:
    \\ 以下哪一项可以帮助减少一组衍生品交易对手方信用风险敞口，除了：
\end{question}

\begin{choices}
    \item Daily mark-to-market of transactions\\
    每日盯市（Daily mark-to-market）
    \item Netting agreements with the counterparty\\
    与交易对手的净额结算协议
    \item Collateral and other credit enhancements\\
    抵押品和其他信用增强
    \item Early termination agreements\\
    提前终止协议
\end{choices}

\begin{correctanswer}
    D
\end{correctanswer}

\begin{explanation}

    \textbf{A选项错误。}
    \begin{itemize}
        \item 每日盯市（Daily mark-to-market）是减少信用风险敞口的有效方法。通过每日重新评估交易价值，可以及时发现风险敞口的变化；盯市机制允许及时调整保证金要求；有助于提前发现潜在的信用问题。
    \end{itemize}

    \textbf{B选项错误。}
    \begin{itemize}
        \item 净额结算协议（Netting agreements）是减少信用风险敞口的重要工具。允许将多个交易的正负敞口相互抵消，显著降低总体信用风险敞口，是场外衍生品市场常用的风险管理工具。
    \end{itemize}

    \textbf{C选项错误。}
    \begin{itemize}
        \item 抵押品和其他信用增强措施是减少信用风险敞口的有效方法。提供额外的信用保护，降低违约时的损失，包括现金抵押、证券抵押等多种形式。
    \end{itemize}

    \textbf{D选项正确。}
    \begin{itemize}
        \item 提前终止协议（Early termination agreements）是减少信用风险敞口效果最差的方法。这些协议通常在交易对手信用质量下降时触发，此时可能已经难以有效执行终止条款，其他交易对手可能已经采取类似行动，导致市场流动性下降，因此效果最差。
    \end{itemize}

\end{explanation}

\checklist{场外衍生品市场的规模比较}


\begin{question}
    Which of the following statements regarding the International Swaps and Derivatives Association (ISDA) Master Agreement is false?
    \\ 哪一项关于国际互换和衍生品协会（ISDA）主协议的陈述是不正确的？
\end{question}

\begin{choices}
    \item The ISDA Master Agreement is a contract between two parties that governs derivative transactions and credit controls\\
    ISDA主协议是双方之间的合同，用于管理衍生品交易和信用控制
    \item An ISDA Master Agreement is prepared for every derivatives transaction that occurs between two parties\\
    ISDA主协议为每笔衍生品交易准备了一份协议
    \item The ISDA Master Agreement forms a single contract under which multiple derivative contracts can be managed\\
    ISDA主协议形成一个单一合同，可以管理多个衍生品合同
    \item The ISDA Master Agreement specifies termination events that each pair of counterparties must accept\\
    ISDA主协议规定了交易双方必须接受的终止事件
\end{choices}

\begin{correctanswer}
    B
\end{correctanswer}

\begin{explanation}

    \textbf{A选项说法恰当。}
    \begin{itemize}
        \item ISDA主协议确实是双方之间的合同，用于管理衍生品交易和信用控制。它规定了交易双方的权利和义务，包含了信用风险控制条款，是场外衍生品交易的基础法律文件。
    \end{itemize}

    \textbf{B选项说法不恰当。}
    \begin{itemize}
        \item ISDA主协议不是为每笔交易单独准备一份协议，而是一个框架协议，可以管理多个衍生品合同。一旦签署，后续交易只需通过确认书（confirmation）确认具体条款。
    \end{itemize}

    \textbf{C选项说法恰当。}
    \begin{itemize}
        \item ISDA主协议确实是一个单一合同，可以管理多个衍生品合同。它提供了标准化的法律框架，允许在同一个协议下进行多笔交易，简化了交易管理流程。
    \end{itemize}

    \textbf{D选项说法恰当。}
    \begin{itemize}
        \item ISDA主协议确实规定了交易双方必须接受的终止事件，包括违约事件和终止事件，并规定了终止后的处理程序。
    \end{itemize}

\end{explanation}

\checklist{ISDA主协议的性质与功能}


\begin{question}
    Which of the following features is least likely a benefit of collateralization?
    \\ 哪一项功能最不可能是一种抵押品的好处？
\end{question}

\begin{choices}
    \item Reduces capital requirements\\
    降低资本要求
    \item Allows for more competitive pricing of counterparty risk\\
    允许对交易对手风险进行更具竞争力的定价
    \item Reduces market, operational, and liquidity risk\\
    降低市场、操作和流动性风险
    \item Reduces credit exposure\\
    降低信用风险敞口
\end{choices}

\begin{correctanswer}
    C
\end{correctanswer}

\begin{explanation}

    \textbf{A选项错误。}
    \begin{itemize}
        \item 抵押品确实可以降低资本要求。通过提供额外的信用保护，降低风险加权资产，从而提高资本效率。
    \end{itemize}

    \textbf{B选项错误。}
    \begin{itemize}
        \item 抵押品确实允许对交易对手风险进行更具竞争力的定价。通过降低信用风险溢价，使交易条件更具吸引力，提高市场竞争力。
    \end{itemize}

    \textbf{C选项正确。}
    \begin{itemize}
        \item 抵押品实际上会增加而不是降低市场、操作和流动性风险。
        \item 市场风险：当抵押品价值下降时，可能导致敞口部分或完全无抵押，增加市场风险敞口。
        \item 操作风险：在获取、管理和处置抵押品时可能遇到法律障碍、技术问题或操作失误，增加操作风险。
        \item 流动性风险：在需要时可能难以以公平市场价格快速出售抵押品，或者抵押品本身可能缺乏流动性，增加流动性风险。
    \end{itemize}

    \textbf{D选项错误。}
    \begin{itemize}
        \item 抵押品确实可以降低信用风险敞口。通过提供额外的信用保护，降低违约时的损失，从而改善信用风险状况。
    \end{itemize}

\end{explanation}
\checklist{抵押品对市场、操作和流动性风险的影响}


\begin{question}
    An investor can gain the exposure and return of an underlying loan by:
    \\ 投资者可以通过以下方式获得标的贷款的风险敞口和收益：
\end{question}

\begin{choices}
    \item Being the credit risk buyer in a total return swap\\
    作为总收益互换中的信用风险买方
    \item Being the credit risk seller in a total return swap\\
    作为总收益互换中的信用风险卖方
    \item Entering into an interest-rate swap with a BB credit\\
    进入一个与BB信用等级的利率互换
    \item Pledging compensating balances with the lender bank
    向贷款银行提供补偿性余额
\end{choices}

\begin{correctanswer}
    A
\end{correctanswer}

\begin{explanation}

    \textbf{A选项正确。}
    \begin{itemize}
        \item 在总收益互换中，信用风险买方（总收益接收方）可以获得标的贷款的风险敞口和收益。信用风险卖方（总收益支付方）将标的资产的总收益支付给买方，买方获得标的贷款的所有收益（包括利息收入和资本利得/损失），并承担标的贷款的信用风险和市场风险。
    \end{itemize}

    \textbf{B选项错误。}
    \begin{itemize}
        \item 作为总收益互换中的信用风险卖方，投资者实际上是转移了标的贷款的风险和收益，而不是获得这些风险和收益。卖方需要支付标的资产的总收益，收取LIBOR加上价差，从而转移了信用风险和市场风险。
    \end{itemize}

    \textbf{C选项错误。}
    \begin{itemize}
        \item 与BB级信用的对手方进行利率互换并不能获得标的贷款的风险敞口和收益。利率互换只涉及利率的交换，不涉及信用风险的转移。
    \end{itemize}

    \textbf{D选项错误。}
    \begin{itemize}
        \item 向贷款银行提供补偿性余额并不能获得标的贷款的风险敞口和收益。补偿性余额只是作为贷款的担保，不涉及风险和收益的转移。
    \end{itemize}

\end{explanation}

\checklist{总收益互换（Total Return Swap）的支付义务}


\begin{question}
    Which of the following statements is not correct regarding total return swaps (TRS)?
    \\ 哪一项关于总收益互换（TRS）的陈述是不正确的？
\end{question}

\begin{choices}
    \item A TRS is designed to mirror the return on an underlying asset like a loan, stock, or even a portfolio of assets\\
    总收益互换（TRS）设计目的是复制标的资产（如贷款、股票或资产组合）的收益
    \item The payer pays any depreciation in the underlying asset to the receiver\\
    支付方支付标的资产的贬值部分给接收方
    \item The payer pays any dividends or interest received to the receiver\\
    支付方支付收到的任何股息或利息给接收方
    \item The receiver is creating a synthetic long position in the underlying\\
    接收方创建一个标的资产的合成多头头寸
\end{choices}

\begin{correctanswer}
    B
\end{correctanswer}
\begin{explanation}

    \textbf{A选项说法恰当。}
    \begin{itemize}
        \item 总收益互换（TRS）的设计目的就是复制标的资产（如贷款、股票或资产组合）的收益，包括利息或股息收入以及资本利得或损失，完全反映标的资产的表现。
    \end{itemize}

    \textbf{B选项说法不恰当。}
    \begin{itemize}
        \item 在总收益互换中，支付方（payer）只负责支付标的资产的升值部分以及相关的股息或利息，而接收方（receiver）需要向支付方支付标的资产的贬值部分。因此，是接收方（而不是支付方）需要支付标的资产的贬值部分给支付方，该说法将责任方搞反了。
    \end{itemize}

    \textbf{C选项说法恰当。}
    \begin{itemize}
        \item 支付方确实需要将收到的任何股息或利息支付给接收方，这是总收益互换的基本机制之一。
    \end{itemize}

    \textbf{D选项说法恰当。}
    \begin{itemize}
        \item 接收方通过总收益互换实际上创建了一个对标的资产的合成多头头寸，因为接收方获得标的资产的所有收益（包括升值和收入），同时也承担所有损失（贬值），相当于持有标的资产。
    \end{itemize}

\end{explanation}

\checklist{总收益互换（TRS）的机制与现金流方向}


\begin{question}
    In pricing a first-to-default credit basket swap, which of the following is true, all else being equal?
    \\ 在定价第一个到违约的信用篮子互换时，哪一项是正确的，所有其他条件不变？
\end{question}

\begin{choices}
    \item The lower the correlation between the assets of the basket, the lower the premium\\
    篮子中资产的相关性越低，保费越低
    \item The lower the correlation between the assets of the basket, the higher the premium\\
    篮子中资产的相关性越低，保费越高
    \item The higher the correlation between the assets of the basket, the higher the premium\\
    篮子中资产的相关性越高，保费越高
    \item The correlation between the assets has no impact in the premium of a first-to-default credit basket swap\\
    篮子中资产的相关性不影响第一个到违约的信用篮子互换的保费
\end{choices}

\begin{correctanswer}
    B
\end{correctanswer}
\begin{explanation}

    \textbf{A选项错误。}
    \begin{itemize}
        \item 这个说法与实际情况相反。在first-to-default信用篮子互换中，资产相关性越低，保费越高，而不是越低。
    \end{itemize}

    \textbf{B选项正确。}
    \begin{itemize}
        \item 在first-to-default信用篮子互换中，资产相关性越低，保费越高。原因如下：
        \begin{itemize}
            \item first-to-default互换在篮子中任何实体首次违约时触发。
            \item 当资产相关性较低时，篮子面临更大的违约风险，因为不同资产可能在不同时间独立违约，增加了至少一个资产违约的概率。
            \item 例如，如果篮子包含来自不同行业的资产，则篮子面临每个行业的独立违约风险，违约概率更高。
            \item 如果篮子只包含来自同一行业的资产（相关性较高），则违约风险较低，因为所有资产可能同时违约或同时不违约，first-to-default触发概率较低。
        \end{itemize}
    \end{itemize}

    \textbf{C选项错误。}
    \begin{itemize}
        \item 这个说法与实际情况相反。在first-to-default信用篮子互换中，资产相关性越高，保费越低，而不是越高。当相关性高时，资产倾向于同时违约或同时不违约，first-to-default触发概率较低。
    \end{itemize}

    \textbf{D选项错误。}
    \begin{itemize}
        \item 资产相关性确实会影响first-to-default信用篮子互换的保费。当资产相关性为负时，first-to-default CDS保费更高；当资产相关性为正且为1时，first-to-default CDS与second-to-default CDS的保费相同。
    \end{itemize}

\end{explanation}

\checklist{first-to-default信用篮子互换（First-to-Default Credit Basket Swap）的保费关系}



\begin{question}
    ABC Company, domiciled in a country with a strong legal system, is applying for a USD 120,000 loan with an annual interest rate of 6\% to be used exclusively on expanding its business and not to repay current debt. The loan will be secured by the company's factory buildings, which have an appraised value of USD 250,000. The company has an annual operating profit of USD 4,000. Which of the following statements relating to the credit risk of this company can most likely be inferred from the given data?
    \\ ABC公司，位于一个法律体系强大的国家，申请了一笔USD 120,000的贷款，年利率为6\%，用于扩大业务，而不是偿还当前债务。贷款将由公司的工厂建筑提供担保，评估价值为USD 250,000。公司年营业利润为USD 4,000。根据给定数据，哪一项关于该公司信用风险的陈述最有可能被推断出来？
\end{question}

\begin{choices}
    \item The company's capacity to pay is low\\
    公司偿付能力很低
    \item The company's willingness to pay is low\\
    公司还款意愿很低
    \item The loan exhibits a high loss given default\\
    贷款在违约时损失很高
    \item The loan's exposure at default will increase\\
    贷款在违约时风险敞口会增加
\end{choices}

\begin{correctanswer}
    A
\end{correctanswer}

\begin{explanation}

    \textbf{A选项正确。}
    \begin{itemize}
        \item 从给定数据可以推断出公司的偿付能力（capacity to pay）很低：
        \begin{itemize}
            \item 贷款金额：\$120,000，年利率：6\%
            \item 年利息支出：
            \[
            \$120,000 \times 6\% = \$7,200
            \]
            \item 公司年营业利润：\$4,000
            \item 营业利润不足以支付利息支出：$\$4,000 < \$7,200$
            \item 更不用说偿还本金，这表明公司的偿付能力很低。
        \end{itemize}
        \item 偿付能力（capacity to pay）是指借款人产生足够现金流来偿还债务的能力，这里营业利润明显不足以覆盖利息支出。
    \end{itemize}

    \textbf{B选项错误。}
    \begin{itemize}
        \item 从给定数据无法推断出公司的还款意愿（willingness to pay）很低。还款意愿是指借款人是否有意愿偿还债务，这通常与道德风险、欺诈行为等有关。题目说明公司申请贷款用于业务扩张，而不是偿还现有债务，这表明公司有良好的还款意愿。公司位于法律体系健全的国家，且提供抵押品，这也表明公司有还款意愿。
    \end{itemize}

    \textbf{C选项错误。}
    \begin{itemize}
        \item 从给定数据无法推断出贷款在违约时会有高损失（LGD）。损失给定违约（Loss Given Default, LGD）取决于回收率，而回收率很大程度上取决于抵押品价值。抵押物（工厂建筑）的评估价值为\$250,000，远高于贷款金额\$120,000。因此，即使发生违约，通过处置抵押品，损失可能较低，而不是高损失。
    \end{itemize}

    \textbf{D选项错误。}
    \begin{itemize}
        \item 从给定数据无法推断出违约时的风险敞口（EAD）会增加。违约时的风险敞口（Exposure at Default, EAD）是指违约发生时对借款人的总敞口。题目中的贷款金额是固定的\$120,000，且贷款用于业务扩张而非偿还现有债务。没有信息表明这个金额会随时间增加，因此无法推断EAD会增加。
    \end{itemize}

\end{explanation}

\checklist{公司的偿付能力分析}


\begin{question}
    Lin Ping is valuing a 1-year credit default swap (CDS) contract which will pay the buyer 80\% of the face value of a bond issued by Xiao Corp, immediately after a default by Xiao. To purchase this CDS, the buyer will pay the CDS spread, which is a percentage of the face value, once at the end of the year. Lin estimates that the risk-neutral default probability for Xiao is 6\% per year. The risk-free rate is 4\% per year. Assuming defaults can only occur halfway through the year and that the accrued premium is paid immediately after a default, what is the estimate for the CDS spread?
    \\ 林平正在估值一个1年期的信用违约互换（CDS）合约，该合约将在Xiao Corp违约后立即支付买家80\%的面值债券。为了购买这个CDS，买家将在年末支付CDS利差，CDS利差是面值的百分比。林平估计Xiao Corp的违约概率为6\%每年。无风险利率为4\%每年。假设违约只能发生在半年后，且应计保费在违约后立即支付，CDS利差的估计是多少？
\end{question}

\begin{choices}
    \item 480 basis points\\
    480个基点
    \item 485 basis points\\
    485个基点
    \item 490 basis points\\
    490个基点
    \item 500 basis points\\
    500个基点
\end{choices}

\begin{correctanswer}
    B
\end{correctanswer}

\begin{explanation}

    \textbf{步骤1：确定基本参数}
    \begin{itemize}
        \item 违约概率（PD）：6\%
        \item 损失给定违约（LGD）：80\%（CDS支付80\%面值，损失为80\%）
        \item 无风险利率（r）：4\%
        \item 违约时间点：0.5年（年中）
    \end{itemize}

    \textbf{步骤2：计算预期损失现值}
    \begin{itemize}
        \item 预期损失 = PD × LGD = 6\% × 80\% = 4.8\%
        \item 损失发生在t = 0.5年，折现到t = 0：
        \[
        \text{损失现值} = \frac{4.8\%}{1 + 0.04 \times 0.5} = \frac{4.8\%}{1.02} = 4.7059\%
        \]
    \end{itemize}

    \textbf{步骤3：计算CDS利差}
    \begin{itemize}
        \item 根据无套利原则，保费现值应等于损失现值。
        \item 保费包括：年末支付（概率94\%）和应计保费（违约时支付，概率6\%）。
        \item 设CDS利差为$s$，则：
        \[
        \frac{0.94s}{1.04} + \frac{0.06s \times 0.5}{1.02} = 4.7059\%
        \]
        \item 求解得：$s \approx 4.85\% = 485\text{ 基点}$。
    \end{itemize}

\end{explanation}

\checklist{CDS利差的计算方法}



\newpage


\section{考点2:对手方风险}
\setcounter{question}{0}


\begin{question}
    Which of the following two transactions increases counterparty credit exposure?

    I. Selling a forward contract to the counterparty

    II. Selling a call option to the counterparty
    \\ 哪一项交易会增加交易对手信用风险敞口？

    I. 出售远期合约给对手方

    II. 出售看涨期权给对手方
\end{question}

\begin{choices}
    \item I only
    \item II only
    \item Both
    \item Neither
\end{choices}

\begin{correctanswer}
    A
\end{correctanswer}

\begin{explanation}
    \textbf{陈述I正确。}
    \begin{itemize}
        \item 出售远期合约给对手方会增加交易对手信用风险敞口。
        \item 在远期合约中，卖方承诺在未来以约定的价格向买方交付标的资产。如果合约在到期时对买方有利（即标的资产价格高于约定价格），买方需要履约支付现金，此时卖方面临对手方信用风险。如果买方违约，卖方可能无法收到应得的付款，从而产生信用损失。因此，远期合约的出售会产生正的信用风险敞口。
    \end{itemize}

    \textbf{陈述II不正确。}
    \begin{itemize}
        \item 出售看涨期权给对手方不会增加交易对手信用风险敞口。
        \item 在期权交易中，卖方在交易开始时已经收取了期权费（premium），这是已收到的现金。卖方的义务是确定的，即在买方行权时以约定价格出售标的资产，这个义务不依赖于交易对手的履约能力。如果买方违约（不行权），卖方保留期权费，不会产生损失。因此，期权卖方不面临对手方信用风险敞口。
    \end{itemize}
\end{explanation}

\checklist{区分不同金融产品对对手方信用风险敞口的影响}

\begin{question}
    Which of the following frictions represents an adverse selection problem?
    \\ 哪一种摩擦代表了逆向选择问题？
\end{question}

\begin{choices}
    \item Investor and mortgagor\\
    投资者和抵押贷款人
    \item Originator and arranger\\
    发起人和承销人
    \item Servicer and rating agency\\
    服务商和评级机构
    \item Servicer and mortgagor\\
    服务商和抵押贷款人
\end{choices}

\begin{correctanswer}
    B
\end{correctanswer}

\begin{explanation}

    \textbf{A选项错误。}
    \begin{itemize}
        \item 投资者和抵押贷款人之间不存在典型的逆向选择问题。投资者通过证券化产品间接持有抵押贷款，主要面临的是信息不对称问题，但不是典型的逆向选择情况。
    \end{itemize}

    \textbf{B选项正确。}
    \begin{itemize}
        \item 发起人（Originator）和承销人（Arranger）之间存在逆向选择问题。发起人拥有关于借款人质量的更多信息，可能保留高质量的抵押贷款，而将低质量的抵押贷款（"柠檬"）出售给承销人。这是典型的逆向选择问题，因为信息不对称导致质量较差的资产更容易被交易。
    \end{itemize}

    \textbf{C选项错误。}
    \begin{itemize}
        \item 服务商和评级机构之间不存在典型的逆向选择问题。服务商负责贷款的管理和收款，评级机构负责评估证券化产品的信用风险，两者之间是服务关系，不是逆向选择。
    \end{itemize}

    \textbf{D选项错误。}
    \begin{itemize}
        \item 服务商和抵押贷款人之间不存在典型的逆向选择问题。服务商负责贷款的管理和收款，抵押贷款人是借款人，两者之间是服务关系，不是逆向选择。
    \end{itemize}

\end{explanation}

\checklist{逆向选择问题}


\begin{question}
    Which of the following methods of mitigating counterparty risk is most likely to generate systemic risk?
    \\ 哪一种方法最有可能产生系统性风险？
\end{question}

\begin{choices}
    \item Netting\\
    净额结算
    \item Collateral\\
    抵押品
    \item Hedging\\
    对冲
    \item Central counterparties\\
    中央对手方
\end{choices}

\begin{correctanswer}
    D
\end{correctanswer}

\begin{explanation}

    \textbf{A选项错误。}
    \begin{itemize}
        \item 净额结算（Netting）不会产生系统性风险。净额结算只是将多个交易的正负敞口相互抵消，风险分散在多个交易对手之间，不会导致风险的集中，因此不会产生系统性风险。
    \end{itemize}

    \textbf{B选项错误。}
    \begin{itemize}
        \item 抵押品（Collateral）不会产生系统性风险。抵押品是双边交易中的信用保护措施，风险分散在多个交易对手之间，不会导致风险的集中，因此不会产生系统性风险。
    \end{itemize}

    \textbf{C选项错误。}
    \begin{itemize}
        \item 对冲（Hedging）不会产生系统性风险。对冲是将风险转移给其他市场参与者，风险分散在多个交易对手之间，不会导致风险的集中，因此不会产生系统性风险。
    \end{itemize}

    \textbf{D选项正确。}
    \begin{itemize}
        \item 中央对手方（Central counterparties）最可能产生系统性风险。中央对手方将原本分散在多个交易对手之间的信用风险集中在少数几个机构中。如果其中一个中央对手方失败或无法履行其义务，可能对整个金融系统造成重大冲击，因为大量交易都依赖于这些中央对手方。这种风险的集中可能导致系统性风险，成为"太大而不能倒"的问题。
    \end{itemize}

\end{explanation}

\checklist{中央对手方（Central Counterparties）的系统性风险}


\begin{question}
    Liz Parker is a junior quantitative analyst who is preparing a report dealing with credit migration. An excerpt of her report contains the following statements:
    \begin{itemize}
    \item I. Future default probability will likely increase over time, especially for periods far into the future.
    \item II. When computing the default probability of a counterparty under a risk-neutral measure, we need to first determine the actual default probability.
    \end{itemize}
    Which of Parker's statements is (are) correct?
    \\ Liz Parker是一名初级量化分析师，正在准备一份关于信用迁移的报告。她的报告摘录包含以下陈述：
    \begin{itemize}
    \item I. 未来违约概率可能会随着时间推移而增加，尤其是对于远期时间点。
    \item II. 在风险中性测度下计算违约概率时，我们需要首先确定实际违约概率。
    \end{itemize}
    Parker的哪一项陈述是正确的？
\end{question}

\begin{choices}
    \item I only
    \item II only
    \item Both I and II
    \item Neither I nor II
\end{choices}

\begin{correctanswer}
    D
\end{correctanswer}

\begin{explanation}
    \textbf{陈述I不正确。}
    \begin{itemize}
        \item 未来的违约概率实际上会随着时间推移而降低，而不是增加。这是因为违约更可能发生在较早的时间点，随着时间推移，如果债务人还没有违约，说明其信用质量可能较好，或者违约已经发生的可能性增加，因此未来违约概率会降低。这是信用迁移中的"幸存者偏差"效应。
    \end{itemize}

    \textbf{陈述II不正确。}
    \begin{itemize}
        \item 在风险中性测度下计算违约概率时，不需要先确定实际违约概率。风险中性测度使用市场隐含的理论概率（从市场价格如CDS利差、债券价格等推导出来），而实际违约概率适用于实际（历史）测度。两者是不同的概念和计算方法，风险中性概率可以直接从市场价格中推导，不需要先计算实际概率。
    \end{itemize}
\end{explanation}

\checklist{未来违约概率和风险中性测度}


\begin{question}
    Which of the following statements regarding counterparty credit risk is most accurate?
    \\ 哪一项关于交易对手信用风险的说法最准确？
\end{question}

\begin{choices}
    \item Counterparty risk is unilateral\\
    交易对手风险是单边的
    \item Over-the-counter (OTC) derivatives contain less counterparty risk than exchange-traded derivatives because the counterparty is known\\
    场外衍生品（OTC）比交易所交易的衍生品风险更小，因为交易对手是已知的
    \item The precise future value of the contract is uncertain, but the counterparties are aware of whether the future value will be positive or negative\\
    合同的确切未来价值不确定，但交易对手知道未来价值是正是负
    \item Counterparty risk is typically associated with counterparty default prior to the settlement rather than default during the settlement process\\
    交易对手风险通常与结算前的违约有关，而不是结算过程中的违约
\end{choices}

\begin{correctanswer}
    D
\end{correctanswer}

\begin{explanation}

    \textbf{A选项错误。}
    \begin{itemize}
        \item 交易对手风险不是单边的，而是双边的。双方都不知道合同的最终价值，双方都不知道是否会盈利或亏损，双方都面临交易对手违约的风险。因此，交易对手风险是双向的，不是单边的。
    \end{itemize}

    \textbf{B选项错误。}
    \begin{itemize}
        \item 场外衍生品的交易对手风险实际上大于交易所交易的衍生品。交易所交易的衍生品的交易对手是交易所，交易所通过保证金制度、每日盯市和清算机制有效降低了交易对手风险。而场外衍生品是双边交易，缺乏这些保护措施，因此交易对手风险更大。
    \end{itemize}

    \textbf{C选项错误。}
    \begin{itemize}
        \item 交易对手不仅不知道合同的确切未来价值，也不知道未来价值是正是负。合同价值取决于未来的市场条件，双方都无法预知未来价值的方向（正或负），这是交易对手风险的本质特征。如果双方能够提前知道价值方向，就不会存在交易对手风险。
    \end{itemize}

    \textbf{D选项正确。}
    \begin{itemize}
        \item 交易对手风险主要与结算前的违约有关，而不是结算过程中的违约。虽然违约可能发生在结算前和结算期间，但交易对手风险主要关注结算前的违约风险，因为此时合同价值可能已经发生重大变化，而对手方在结算前违约会导致无法收取应得的收益或避免损失。结算过程中的违约虽然也可能发生，但不是交易对手风险的主要关注点。
    \end{itemize}

\end{explanation}

\checklist{交易对手风险与贷款风险的关系}


\begin{question}
    Which of the following statements regarding lending and counterparty risk is true?
    \\ 哪一项关于贷款和交易对手风险的说法是正确的？
\end{question}

\begin{choices}
    \item Counterparty risk is defined as the possibility that an obligator will default on an outstanding loan\\
    交易对手风险定义为借款人可能违约的可能性
    \item Lending risk is defined as the possibility that either party to a derivative transaction will fail to meet its obligations\\
    贷款风险定义为衍生品交易中任何一方可能违约的可能性
    \item Counterparty risk is very different from and more complex than lending risk\\
    交易对手风险与贷款风险有很大区别，且更复杂
    \item Counterparty risk and lending risk are the same\\
    交易对手风险和贷款风险是相同的
\end{choices}

\begin{correctanswer}
    C
\end{correctanswer}

\begin{explanation}

    \textbf{A选项错误。}
    \begin{itemize}
        \item 这个定义实际上是贷款风险的定义，而不是交易对手风险的定义。贷款风险是指借款人（obligator）可能违约的可能性，而交易对手风险是指在衍生品交易中任何一方可能违约的可能性。这个选项混淆了两种风险的定义。
    \end{itemize}

    \textbf{B选项错误。}
    \begin{itemize}
        \item 这个定义实际上是交易对手风险的定义，而不是贷款风险的定义。交易对手风险是指衍生品交易中任何一方可能违约的可能性，而贷款风险是指借款人可能违约的可能性。这个选项混淆了两种风险的定义。
    \end{itemize}

    \textbf{C选项正确。}
    \begin{itemize}
        \item 交易对手风险与贷款风险有很大区别，且更复杂。交易对手风险是双向的，双方都可能违约；涉及未来不确定的现金流，合同价值取决于市场条件；需要考虑市场风险和信用风险的相互作用；交易对手风险的管理更复杂，因为需要同时考虑市场风险和信用风险。相比之下，贷款风险是单向的，只关注借款人的违约，现金流相对确定。
    \end{itemize}

    \textbf{D选项错误。}
    \begin{itemize}
        \item 交易对手风险和贷款风险是不同的。贷款风险是单向的，只关注借款人的违约，而交易对手风险是双向的，双方都可能违约。两种风险的特征和管理方法都不同，因此不能视为相同。
    \end{itemize}

\end{explanation}

\checklist{交易对手风险与贷款风险的区别}


\begin{question}
    Which of the following statements is not a motivation for pricing counterparty risk?
    \\ 哪一项不是交易对手风险定价的动机？
\end{question}

\begin{choices}
    \item Accurate pricing should only account for the cost of the trade\\
    准确的定价只应考虑交易成本
    \item Counterparty risk pricing should account for risk mitigants\\
    交易对手风险定价应考虑风险缓释措施
    \item Best practices organize pricing responsibilities in the organization\\
    最佳实践组织定价责任在组织中
    \item Pricing bilateral derivatives contracts\\
    为双边衍生品合约定价
\end{choices}

\begin{correctanswer}
    A
\end{correctanswer}

\begin{explanation}

    \textbf{A选项正确。}
    \begin{itemize}
        \item 这个说法不是交易对手风险定价的动机。准确的定价不仅要考虑交易成本，还要考虑交易对手风险成本、信用风险溢价等因素。只考虑交易成本是不完整的，因为这会忽略交易对手风险带来的额外成本，导致定价不准确。因此，这个说法不是交易对手风险定价的动机。
    \end{itemize}

    \textbf{B选项错误。}
    \begin{itemize}
        \item 考虑风险缓释措施是交易对手风险定价的重要动机。风险缓释措施（如抵押品、净额结算、担保等）可以降低交易对手风险，需要在定价中反映这些措施的效果，以准确评估和管理风险。这是风险定价的重要组成部分。
    \end{itemize}

    \textbf{C选项错误。}
    \begin{itemize}
        \item 在组织中明确定价责任是交易对手风险定价的动机。确保风险定价的准确性和一致性，明确各部门（如交易部门、风险管理部门）的职责，这是良好的风险管理实践，有助于建立有效的风险管理框架。
    \end{itemize}

    \textbf{D选项错误。}
    \begin{itemize}
        \item 为双边衍生品合约定价是交易对手风险定价的动机。双边衍生品合约（如场外衍生品）涉及交易对手风险，需要在定价中反映这种风险，以准确评估交易的价值和风险，这是风险定价的重要应用场景。
    \end{itemize}

\end{explanation}

\checklist{交易对手风险定价的动机}


\begin{question}
    A risk manager at a large regional bank asks a group of junior analysts to study and report on the economic costs of OTC derivatives traded by the bank in relation to the bank's counterparty risk and funding risk. The manager instructs the analysts to use a hypothetical threshold that is higher for the counterparty than for the bank in assessing the bank's funding and collateral positions. Which of the following observations, from the perspective of the regional bank, would the analysts be correct to make?
    \\ 一家大型区域银行的风险经理要求一群初级分析师研究并报告银行交易对手风险和融资风险相关的OTC衍生品的经济成本。经理指示分析师在评估银行的融资和抵押品时，使用比银行更高的对手方阈值。以下哪一项观察，从区域银行的角度来看，分析师是正确的？
\end{question}

\begin{choices}
    \item The bank has a funding cost but faces no counterparty risk when a derivative portfolio is in-the-money\\
    银行在衍生品组合处于实值状态时面临融资成本，但不会面临交易对手风险
    \item The bank faces no counterparty risk and no funding cost when a derivative portfolio is out-of-the-money\\
    银行在衍生品组合处于虚值状态时不会面临交易对手风险，也不会面临融资成本
    \item The bank has a potential funding benefit when a derivative portfolio is in-the-money but not when it is out-of-the-money\\
    银行在衍生品组合处于实值状态时面临融资收益，但在虚值状态时没有融资收益
    \item The bank faces counterparty risk whether a derivative portfolio is in-the-money or out-of-the-money\\
    银行在衍生品组合处于实值状态或虚值状态时都面临交易对手风险
\end{choices}

\begin{correctanswer}
    B
\end{correctanswer}

\begin{explanation}

    \textbf{A选项错误。}
    \begin{itemize}
        \item 当衍生品组合处于实值状态（in-the-money）时，银行确实面临融资成本。因为银行应收钱（合约价值为正），而对手方的阈值更高意味着银行需要等待更长时间才能收到抵押品，从而产生融资成本。但银行仍然面临交易对手风险，因为如果对手方违约，银行将无法收到应得的款项，因此这个观察是错误的。
    \end{itemize}

    \textbf{B选项正确。}
    \begin{itemize}
        \item 当衍生品组合处于虚值状态（out-of-the-money）时，银行不面临交易对手风险，也不面临融资成本。因为银行应付钱（合约价值为负），对手方应收钱，但对手方的阈值更高意味着对手方需要等待更长时间才能收到抵押品。如果对手方违约，银行不会遭受损失（因为银行是付款方），因此没有交易对手风险。同时，银行不需要提供抵押品（因为银行的头寸为负且对手方阈值更高），因此没有融资成本。这是正确的观察。
    \end{itemize}

    \textbf{C选项错误。}
    \begin{itemize}
        \item 当衍生品组合处于实值状态时，银行面临融资成本，而不是融资收益。因为银行应收钱，需要等待更长时间才能收到抵押品，从而产生融资成本。当衍生品组合处于虚值状态时，银行不需要提供抵押品，因此也没有融资收益。这个观察是错误的。
    \end{itemize}

    \textbf{D选项错误。}
    \begin{itemize}
        \item 银行面临的交易对手风险取决于衍生品组合的价值状态。当组合处于虚值状态时，银行不面临交易对手风险（因为银行是付款方，对手方违约不会造成损失）。只有当组合处于实值状态时，银行才面临交易对手风险（因为银行是收款方，对手方违约会造成损失）。因此，这个观察是错误的。
    \end{itemize}

\end{explanation}

\checklist{衍生品组合价值状态与交易对手风险和融资成本的关系}






\section{考点3:净额结算及相关方面}
\setcounter{question}{0}

\begin{question}
    Company ABC has great deal of transactions with counterparty XYZ. ABC has some contracts with XYZ covered under netting agreement A, has some contracts with XYZ covered under netting agreement B, and some contracts with XYZ with no netting agreement. With the following data, what is ABC's current counterparty credit exposure to XYZ?
    \\ ABC公司与交易对手XYZ有很多交易。ABC有一些合同在净额协议A下覆盖，有一些合同在净额协议B下覆盖，还有一些合同在XYZ下没有净额协议。根据以下数据，ABC对XYZ的当前交易对手信用风险敞口是多少？
\end{question}

\begin{choices}
    \item USD 8,500
    \item USD 7,200
    \item USD 5,800
    \item USD 10,200
\end{choices}

\begin{correctanswer}
    D
\end{correctanswer}

\begin{explanation}

    \textbf{步骤1：计算净额协议A下的敞口}
    \begin{itemize}
        \item 净额协议A下的净敞口：
        \[
        \text{净敞口} = 2,500 - 3,800 - 2,200 = -3,500
        \]
        \item 信用风险敞口取正值（如果为负则取0）：
        \[
        \text{敞口}_A = \max(-3,500, 0) = 0
        \]
    \end{itemize}

    \textbf{步骤2：计算净额协议B下的敞口}
    \begin{itemize}
        \item 净额协议B下的净敞口：
        \[
        \text{净敞口} = 6,500 - 800 - 2,800 + 4,500 = 7,400
        \]
        \item 信用风险敞口：
        \[
        \text{敞口}_B = \max(7,400, 0) = 7,400
        \]
    \end{itemize}

    \textbf{步骤3：计算无净额协议下的敞口}
    \begin{itemize}
        \item 无净额协议下的交易不能与其他交易抵消，因此敞口为：$\$2,800$。
    \end{itemize}

    \textbf{步骤4：计算总敞口}
    \begin{itemize}
        \item 总交易对手信用风险敞口：
        \[
        \text{总敞口} = 0 + 7,400 + 2,800 = 10,200
        \]
    \end{itemize} 

  
\end{explanation}

\checklist{交易对手信用风险敞口的计算方法}


\begin{question}
    Bank One has made a \$250 million loan to a software company at a fixed rate of 10 percent. The bank wants to hedge its exposure by entering into a Total Return Swap with a counterparty, Interloan Co., in which Bank One promises to pay the interest on the loan plus the change in the market value of the loan in exchange for LIBOR plus 50 basis points. If after one year, the market value of the loan has decreased by 4 percent and LIBOR is 9 percent, what will be the net obligation of Bank One?
    \\ Bank One向一家软件公司发放了\$250百万的贷款，固定利率为10\%。银行希望通过与交易对手Interloan Co.签订总收益互换合约来对冲其风险，Bank One承诺支付贷款利息加上贷款市场价值的变动，而Interloan Co.承诺支付LIBOR加上50个基点。一年后，贷款的市场价值下降了4\%，LIBOR为9\%，Bank One的净义务是多少？
\end{question}

\begin{choices}
    \item Net receipt of \$6.0 million\\
    净收入\$6.0百万
    \item Net payment of \$6.0 million\\
    净支付\$6.0百万
    \item Net receipt of \$6.5 million\\
    净收入\$6.5百万
    \item Net payment of \$6.5 million\\
    净支付\$6.5百万
\end{choices}

\begin{correctanswer}
    A
\end{correctanswer}

\begin{explanation}
    \textbf{A选项正确。}
    \begin{itemize}
        \item 银行应收金额 = \$250 million × (9\% + 0.5\%) = \$23.75 million
        \item 银行应付金额 = \$250 million × 10\% + \$250 million × (-4\%) = \$15 million
        \item 净收入 = \$23.75 million - \$15 million = \$6.0 million
    \end{itemize}

\end{explanation}

\checklist{总收益互换（Total Return Swap）的净义务计算}

\begin{question}
    A diversified portfolio of OTC derivatives with a single counterparty currently has a net mark-to-market of USD 25,000,000 and a gross mark-to-market (the sum of the value of all positive value positions minus the value of all negative value positions) of USD 90,000,000. Assuming there are no netting agreements in place with the counterparty, determine the current credit exposure to the counterparty.
    \\ 一个单一交易对手的OTC衍生品组合目前有USD 25百万的净Mark-to-Market，总Mark-to-Market（所有正价值头寸的价值之和减去所有负价值头寸的价值之和）为USD 90百万。假设与交易对手没有净额协议，确定当前对交易对手的信用风险敞口。
\end{question}

\begin{choices}
    \item Less than or equal to USD 15,000,000\\
    小于或等于USD 15百万
    \item Greater than USD 15,000,000 but less than or equal to USD 45,000,000\\
    大于USD 15百万但小于或等于USD 45百万
    \item Greater than USD 45,000,000 but less than USD 65,000,000\\
    大于USD 45百万但小于USD 65百万
    \item Greater than USD 65,000,000\\
    大于USD 65百万
\end{choices}

\begin{correctanswer}
    C
\end{correctanswer}

\begin{explanation}


    \textbf{C选项正确。}
    \begin{itemize}
        \item 设正价值头寸为A，负价值头寸为B
        \item 净值 = A - B = \$25 million
        \item 总值 = A + B = \$90 million
        \item 解得：A = \$57.5 million
        \item 因此，当前信用风险敞口为\$57.5 million
    \end{itemize}

\end{explanation}

\checklist{无净额协议时，信用风险敞口的计算方法}


\begin{question}
    Riggs Resources, LLC, (Riggs) is a commodity trading firm. Riggs has numerous trades outstanding with several counterparties; however, it is concerned with presettlement risk. In order to reduce presettlement risk (the risk that Riggs's counterparties would default before settlement), it would be most beneficial for Riggs to
    \\ Riggs Resources, LLC（Riggs）是一家商品交易公司。Riggs与多个交易对手有众多未结算交易，但担心结算前风险。为了降低结算前风险（Riggs的交易对手在结算前违约的风险），对Riggs最有利的是
\end{question}

\begin{choices}
    \item Have payment netting\\
    有支付净额结算
    \item Have close-out netting\\
    有终止净额结算
    \item Analyze potential losses as the sum of exposures\\
    将潜在损失作为敞口之和进行分析
    \item Have netting but not set-off\\
    有净额结算但不是抵销
\end{choices}

\begin{correctanswer}
    B
\end{correctanswer}

\begin{explanation}

    \textbf{A选项错误。}
    \begin{itemize}
        \item 支付净额结算（payment netting）不能有效降低结算前风险。支付净额结算主要降低结算风险和操作风险，它涉及将同一天到期的多个支付进行净额结算，但不涉及交易对手违约的处理，因此不能有效降低结算前风险（presettlement risk）。
    \end{itemize}

    \textbf{B选项正确。}
    \begin{itemize}
        \item 终止净额结算（close-out netting）最能有效降低结算前风险。在交易对手违约时，终止净额结算允许终止所有合同并进行净额结算，将多个交易的正负敞口相互抵消，显著降低结算前风险敞口。这是降低结算前风险最有效的方法。
    \end{itemize}

    \textbf{C选项错误。}
    \begin{itemize}
        \item 将敞口简单相加分析潜在损失不能降低结算前风险。这种方法忽略了净额结算的效果，会高估实际风险敞口，不能有效降低风险，只是风险度量方法，而不是风险降低措施。
    \end{itemize}

    \textbf{D选项错误。}
    \begin{itemize}
        \item 净额结算和抵销（set-off）实际上是同义词，都指将多个交易的正负敞口相互抵消。不能同时有净额结算而没有抵销，这个选项本身是矛盾的，因此不正确。
    \end{itemize}

\end{explanation}

\checklist{终止净额结算（Close-out Netting）的优点}


\begin{question}
    What does netting permit a firm to do?
    \\ 净额结算允许公司做什么？
\end{question}

\begin{choices}
    \item Pay only the losses due to unfavorable market prices\\
    只支付由于不利市场价格造成的损失
    \item Net all of its gains against all of its losses\\
    将其所有收益与其所有损失相互抵消
    \item Net its transaction with a given counterparty against each other\\
    将其与特定交易对手的交易相互抵消
    \item Pay the losses only if the counterpart defaults\\
    只有在交易对手违约时才支付损失
\end{choices}

\begin{correctanswer}
    C
\end{correctanswer}

\begin{explanation}

    \textbf{A选项错误。}
    \begin{itemize}
        \item 净额结算不是只支付不利市场价格造成的损失。净额结算涉及所有交易的正负敞口相互抵消，不仅限于市场风险造成的损失，而是将公司与特定交易对手之间的所有应收和应付款项进行净额结算，这个描述不准确。
    \end{itemize}

    \textbf{B选项错误。}
    \begin{itemize}
        \item 净额结算不是将所有收益和损失相互抵消。净额结算只适用于与同一交易对手的交易，不能跨交易对手进行净额结算。不同交易对手之间的交易不能相互抵消，因此这个描述不准确。
    \end{itemize}

    \textbf{C选项正确。}
    \begin{itemize}
        \item 净额结算允许公司将与特定交易对手的交易相互抵消。净额结算将应收交易对手的款项与应付交易对手的款项相抵消，减少实际支付金额，从而降低信用风险敞口。这是净额结算的核心功能。
    \end{itemize}

    \textbf{D选项错误。}
    \begin{itemize}
        \item 净额结算不是只在交易对手违约时才支付损失。净额结算是正常的结算机制，在正常的交易结算过程中使用，不依赖于交易对手是否违约。无论是正常结算还是违约情况下的终止净额结算，都是将正负敞口相互抵消，这个描述不准确。
    \end{itemize}

\end{explanation}
\checklist{净额结算（Netting）的功能}


\begin{question}
    Ondine Financial Inc. (Ondine) uses a variety of techniques to manage counterparty risk. It has entered into an interest rate swap with Scarbo Inc. (Scarbo). Currently, Ondine's position in the swap has a \$1.5 million mark-to-market value. Based on the information provided, which of the following credit risk mitigation techniques would be most advantageous to Ondine if Scarbo defaults?
    \\ Ondine Financial Inc.（Ondine）使用多种技术管理交易对手风险。它与Scarbo Inc.（Scarbo）签订了利率互换合约。目前，Ondine在互换中的头寸有\$1.5百万的Mark-to-Market价值。根据提供的信息，如果Scarbo违约，哪一种信用风险缓释技术对Ondine最有利？
\end{question}

\begin{choices}
    \item Close-out\\
    终止
    \item Collateralization\\
    抵押品
    \item Netting\\
    净额结算
    \item Walkaway\\
    撕毁
\end{choices}

\begin{correctanswer}
    D
\end{correctanswer}

\begin{explanation}

    \textbf{A选项错误。}
    \begin{itemize}
        \item 终止（Close-out）在这种情况下对Ondine不利。如果Ondine的合约价值为负（即Ondine需要支付），终止后需要支付净额，不能减少损失，反而会使损失确定化。
    \end{itemize}

    \textbf{B选项错误。}
    \begin{itemize}
        \item 抵押品（Collateralization）在这种情况下不相关或不能有效减少损失。如果Ondine是净支付方（合约价值为负），Ondine不需要向对手方提供抵押品，因此抵押品缓释技术在这种情况下不能帮助Ondine减少损失。
    \end{itemize}

    \textbf{C选项错误。}
    \begin{itemize}
        \item 净额结算（Netting）在这种情况下对Ondine不利。如果Ondine需要支付净额，净额结算只是将多个交易的正负敞口相互抵消，但Ondine仍然需要支付净额，不能减少损失。
    \end{itemize}

    \textbf{D选项正确。}
    \begin{itemize}
        \item "撕毁"条款（Walkaway）在这种情况下对Ondine最有利。如果Ondine的合约价值为负（即Ondine需要支付），walkaway条款允许Ondine在对手方违约时取消交易而不需要支付任何款项，从而可以最大限度地减少或避免损失。这是唯一能够在对手方违约时帮助Ondine避免支付义务的缓释技术。
    \end{itemize}

\end{explanation}

\checklist{信用风险缓释技术在不同合约价值情况下的适用性与效果}



\begin{question}
    The following scenarios illustrate exposures for two trades with and without the impact of netting. Given the below expected positive exposures (EPE), what is the corresponding netting factor? Assume that the average correlation between exposures is 0.65.
    \\ 以下场景说明了有和没有净额结算影响的两种交易的风险敞口。根据以下预期正敞口（EPE），对应的净额因子是多少？假设敞口之间的平均相关系数为0.65。
\end{question}

\begin{choices}
    \item 0.65
    \item 0.75
    \item 0.85
    \item 1.35
\end{choices}

\begin{correctanswer}
    C
\end{correctanswer}

\begin{explanation}
   
    \textbf{步骤1：确定参数}
    \begin{itemize}
        \item 敞口数量：$n = 5$
        \item 平均相关系数：$\overline{\rho} = 0.65$
    \end{itemize}

    \textbf{步骤2：计算净额因子}
    \begin{itemize}
        \item 净额因子计算公式：
        \[
        \text{净额因子} = \frac{\sqrt{n + n(n-1)\overline{\rho}}}{n}
        \]
        \item 代入数值计算：
        \[
        \begin{aligned}
        \text{净额因子} &= \frac{\sqrt{5 + 5(5-1) \times 0.65}}{5} \\
        &= \frac{\sqrt{5 + 5 \times 4 \times 0.65}}{5} \\
        &= \frac{\sqrt{5 + 13}}{5} \\
        &= \frac{\sqrt{18}}{5} \\
        &= \frac{4.24}{5} \\
        &= 0.85
        \end{aligned}
        \]
    \end{itemize}


 
\end{explanation}

\checklist{净额因子计算}


\begin{question}
    A firm has entered into a USD 25 million total return swap on the NASDAQ 100 Index as the index payer with ABC Corporation, which will pay 1-year LIBOR +3.0\%. The contract will last 1 year, and cash flows will be exchanged annually. Suppose the NASDAQ 100 Index is currently at 3,200 and LIBOR is 1.50\%. The firm conducts a stress test on this total return swap using the following scenario: NASDAQ 100 in 1 year: 4,000 LIBOR in 1 year: 0.75\%.For this scenario, what is the firm's net cash flow in year 1?
    \\ 一家公司与ABC Corporation签订了USD 25百万的纳斯达克100指数总收益互换合约，作为指数支付方。ABC Corporation将支付1年LIBOR +3.0\%。合约为期1年，每年交换现金流。假设纳斯达克100指数目前为3,200，LIBOR为1.50\%。公司对这份总收益互换合约进行压力测试，使用以下场景：1年后纳斯达克100指数为4,000，LIBOR为0.75\%。对于这个场景，公司第1年的净现金流是多少？
\end{question}

\begin{choices}
    \item A net cash outflow of USD 5.50 million\\
    净现金流流出USD 5.50百万
    \item A net cash outflow of USD 5.25 million\\
    净现金流流出USD 5.25百万
    \item A net cash outflow of USD 5.00 million\\
    净现金流流出USD 5.00百万
    \item A net cash outflow of USD 5.75 million\\
    净现金流流出USD 5.75百万
\end{choices}

\begin{correctanswer}
    B
\end{correctanswer}

\begin{explanation}
\textbf{步骤1：计算指数支付金额}
        \begin{itemize}
            \item 指数上涨幅度：
            \[
            \frac{4,000}{3,200} - 1 = 0.25 = 25\%
            \]
            \item 公司需要支付的指数收益：
            \[
            25\% \times \$25\text{ million} = \$6.25\text{ million}
            \]
        \end{itemize}

\textbf{步骤2：计算公司收到的LIBOR支付金额}
        \begin{itemize}
            \item 当前LIBOR = 1.50\%
            \item 利差 = 3.0\%
            \item 总利率 = 1.50\% + 3.0\% = 4.50\%
            \item ABC公司支付的金额：
            \[
            4.50\% \times \$25\text{ million} = \$1.00\text{ million}
            \]
        \end{itemize}

\textbf{步骤3：计算净现金流}
        \begin{itemize}
            \item 净现金流 = 收到的金额 - 支付的金额：
            \[
            \$1.00\text{ million} - \$6.25\text{ million} = -\$5.25\text{ million}
            \]
            \item 负号表示净现金流出。
        \end{itemize}

\end{explanation}

\checklist{总收益互换（Total Return Swap）的现金流计算}




\section{考点4:保证金（抵押品）和结算}
\setcounter{question}{0}


\begin{question}
    Bank A trades a 10-year interest rate swap with Bank B. Bank A is concerned about the 1-year exposure since Bank B has a poor credit rating. Which of the following measures will exacerbate Bank A's credit exposure to Bank B?
    \begin{itemize}
    \item I. Application of CSA
    \item II. Execute the swap as a reset swap wherein the swap will be marked to market every six months
    \item III. Application of break clause
    \item IV. Decrease the frequency of coupon payments
    \end{itemize}
     Bank A与Bank B签订了10年利率互换合约。Bank A担心1年的信用风险敞口，因为Bank B的信用评级较差。哪一项措施会加剧Bank A对Bank B的信用风险敞口？
     \begin{itemize}
    \item I. 应用信用支持协议（CSA）
    \item II. 执行重置互换，每6个月进行盯市
    \item III. 应用终止条款
    \item IV. 减少息票支付频率
    \end{itemize}
\end{question}

\begin{choices}
    \item I only
    \item IV only
    \item I, II, III and IV
    \item I, II and III
\end{choices}

\begin{correctanswer}
    B
\end{correctanswer}

\begin{explanation}
    \textbf{陈述I不正确。}
    \begin{itemize}
        \item 信用支持协议（CSA）不会增加信用风险敞口，而是会降低敞口。CSA要求提供抵押品，可以降低信用风险敞口，是风险缓释措施。
    \end{itemize}

    \textbf{陈述II不正确。}
    \begin{itemize}
        \item 执行重置互换（每6个月进行盯市）不会增加信用风险敞口，而是会降低敞口。重置互换通过定期重置合约价值，可以限制敞口的累积，减少信用风险敞口。
    \end{itemize}

    \textbf{陈述III不正确。}
    \begin{itemize}
        \item 应用终止条款（break clause）不会增加信用风险敞口，而是会降低敞口。终止条款允许在特定条件下提前终止交易，从而限制敞口的累积，减少信用风险敞口。
    \end{itemize}

    \textbf{陈述IV正确。}
    \begin{itemize}
        \item 减少息票支付频率会增加信用风险敞口。支付间隔时间越长，市场利率变动的风险越大，敞口累积的时间更长，增加了交易对手违约时的潜在损失，因此会加剧信用风险敞口。
    \end{itemize}
\end{explanation}

\checklist{减少息票支付频率对信用风险敞口的影响}


\begin{question}
    Which of the following statements is least accurate regarding a credit support annex (CSA) and/or an ISDA Master Agreement?
    \\ 哪一项关于信用支持附件（CSA）和/或ISDA主协议的说法最不准确？
\end{question}

\begin{choices}
    \item ISDA Master Agreements help standardize collateral management\\
    ISDA主协议有助于标准化抵押品管理
    \item CSAs must define all collateralization parameters in order to work as intended\\
    CSA必须定义所有抵押品参数才能按预期工作
    \item Compared to the ISDA Master Agreement, CSAs were first to establish collateral standards\\
    与ISDA主协议相比，CSA是第一个建立抵押品标准的
    \item CSAs are incorporated into an ISDA Master Agreement\\
    CSA被纳入ISDA主协议
\end{choices}

\begin{correctanswer}
    C
\end{correctanswer}

\begin{explanation}

    \textbf{A选项说法恰当。}
    \begin{itemize}
        \item ISDA主协议确实有助于标准化抵押品管理，提供了标准化的法律框架，规定了抵押品管理的基本规则，促进了市场实践的统一。
    \end{itemize}

    \textbf{B选项说法恰当。}
    \begin{itemize}
        \item 信用支持附件（CSA）必须定义所有抵押品参数才能按预期工作，包括抵押品资格、利率支付、转移的时间和机制、抵押品计算、抵押品证券的折扣、抵押品替换、估值时间和方法、抵押品再利用、争议处理等。
    \end{itemize}

    \textbf{C选项说法不恰当。}
    \begin{itemize}
        \item CSA不是最先建立抵押品标准的，ISDA主协议先于CSA建立。CSA是作为ISDA主协议的补充而发展起来的，用于规定抵押品管理的具体细节。因此，相比ISDA主协议，CSA并不是最先建立抵押品标准的，这个说法与历史发展不符，是最不准确的。
    \end{itemize}

    \textbf{D选项说法恰当。}
    \begin{itemize}
        \item CSA确实是ISDA主协议的组成部分，是ISDA主协议的附件，用于规定抵押品管理的具体细节，可以选择单边或双边CSA。
    \end{itemize}

\end{explanation}

    \checklist{CSA（Credit Support Annex）的性质与功能}


\begin{question}
    When dealing with a hedge fund, a bank would most likely negotiate:
    \\ 在与对冲基金打交道时，银行最有可能谈判：
\end{question}

\begin{choices}
    \item A one-way agreement in the bank's favor given the bank's stronger credit rating\\
    银行信用评级更强时，银行最有可能谈判一个对自己有利的单向协议（one-way agreement）。
    \item A one-way agreement in the bank's favor agreeing to post collateral to the hedge fund\\
    银行信用评级更强时，银行最有可能谈判一个对自己有利的单向协议（one-way agreement），同意向对冲基金提供抵押品。
    \item A two-way agreement given the relatively small difference in credit quality between the two entities\\
    银行信用评级更强时，银行最有可能谈判一个对自己有利的双向协议（two-way agreement），同意向对冲基金提供抵押品。
    \item A two-way agreement where both parties agree to post collateral\\
    银行信用评级更强时，银行最有可能谈判一个对自己有利的双向协议（two-way agreement），同意向对冲基金提供抵押品。
\end{choices}

\begin{correctanswer}
    A
\end{correctanswer}

\begin{explanation}

    \textbf{A选项正确。}
    \begin{itemize}
        \item 银行很可能会谈判一个对自己有利的单向协议（one-way agreement）。银行的信用评级通常高于对冲基金，当交易双方信用质量差异较大时，适合签订单向协议。在单向协议中，信用质量较差的一方（对冲基金）需要向信用质量较好的一方（银行）提供抵押品，以降低银行的信用风险敞口。
    \end{itemize}

    \textbf{B选项错误。}
    \begin{itemize}
        \item 银行信用质量更好，不应该向对冲基金提供抵押品。这种安排会增加银行的风险，不符合风险管理的原则。在单向协议中，应该是信用质量较差的一方（对冲基金）向信用质量较好的一方（银行）提供抵押品。
    \end{itemize}

    \textbf{C选项错误。}
    \begin{itemize}
        \item 银行和对冲基金的信用质量通常有显著差异，银行通常具有更高的信用评级。双向协议适用于信用质量相近的交易对手，因此不适合银行和对冲基金之间的情况。
    \end{itemize}

    \textbf{D选项错误。}
    \begin{itemize}
        \item 双向协议适用于信用质量相近的交易对手，双方都同意提供抵押品。但银行和对冲基金的信用质量差异较大，因此这种安排不符合风险管理的最佳实践。银行应该要求对冲基金提供抵押品，而不是双方都提供。
    \end{itemize}

\end{explanation}

\checklist{银行家悖论（Banker's Paradox）的本质与影响}


\begin{question}
    Leverage, Inc., an investment bank, has numerous credit default swaps with XYZ Corp. Leverage has established a break clause with XYZ Corp. to reduce risk. The break clause is trigger-based and may be exercised once the trigger is satisfied. The CEO of Leverage is concerned about a banker's paradox. Which of the following statements best describe the CEO's concern?
    \\ Leverage, Inc.是一家投资银行，与XYZ Corp.签订了多个信用违约互换合约。Leverage与XYZ Corp.签订了中断条款（break clause）以降低风险。中断条款是触发式（trigger-based）的，一旦触发条件满足，就可以行使。Leverage的CEO担心银行家悖论。哪一项陈述最准确地描述了CEO的担忧？
\end{question}

\begin{choices}
    \item To be effective, the break clause option should not be used too early\\
    为了有效，中断条款选项不应该过早使用
    \item The weak firm often recovers after the use of the break clause\\
    弱公司通常在使用了中断条款后恢复
    \item The break clause option is used too late, and the weak firm gets weaker\\
    中断条款选项使用太晚，弱公司变得更弱
    \item The break clause option is used too early, and relations with the counterparty suffer\\
    中断条款选项使用太早，与交易对手的关系受损
\end{choices}

\begin{correctanswer}
    C
\end{correctanswer}

\begin{explanation}

    \textbf{A选项错误。}
    \begin{itemize}
        \item 这个说法与银行家悖论相反。中断条款应该在交易对手信用质量显著下降之前使用，过早使用不是问题所在，主要问题是使用太晚。这个选项误解了银行家悖论的本质。
    \end{itemize}

    \textbf{B选项错误。}
    \begin{itemize}
        \item 这个说法不准确。交易对手在信用质量显著下降后通常不会恢复，这个选项低估了信用风险，不符合实际情况。银行家悖论的核心不是关于交易对手是否会恢复，而是关于中断条款使用的时机问题。
    \end{itemize}

    \textbf{C选项正确。}
    \begin{itemize}
        \item 这个说法准确描述了银行家悖论。中断条款应该在交易对手信用质量显著下降之前使用，但机构通常为了避免损害与交易对手的关系而推迟使用中断条款，导致中断条款使用太晚，此时交易对手的信用质量已经进一步恶化，无法有效降低风险敞口。这就是银行家悖论：最佳的避险时机（信用质量下降初期）与保护客户关系的需要之间存在矛盾。
    \end{itemize}

    \textbf{D选项错误。}
    \begin{itemize}
        \item 这个说法与银行家悖论相反。过早使用不是主要问题，主要问题是使用太晚。银行家悖论的核心是机构倾向于推迟使用中断条款，导致使用时已经太晚，这个选项误解了银行家悖论的本质。
    \end{itemize}

\end{explanation}

\checklist{银行家悖论（Banker's Paradox）的本质与影响}

\begin{question}
    A financial institution has a two-way collateral support annex (CSA) with a counterparty covering a portfolio valued at JPY 500 million. The margining terms of the collateralized portfolio include a threshold of JPY 200 million, a minimum transfer amount of JPY 40 million, and a margin period of risk of 12 days. Which of the following is correct regarding the size of collateral in mitigating the counterparty risk of the portfolio?
    \\ 一家金融机构与交易对手签订了双向抵押品支持附件（CSA），覆盖了一个价值500百万日元的投资组合。抵押品覆盖的投资组合的保证金条款包括一个门槛值为200百万日元，最小转移金额为40百万日元，保证金风险期为12天。哪一项关于抵押品在降低投资组合交易对手风险中的作用的说法是正确的？
\end{question}

\begin{choices}
    \item A lower threshold value is equivalent to a larger portion of exposure protected by collateral\\
    较低的门槛值相当于更大比例的敞口受到抵押品保护
    \item A shorter margin period of risk is equivalent to a smaller portion of exposure protected by collateral\\
    较短的保证金风险期相当于更小比例的敞口受到抵押品保护
    \item A lower independent amount is equivalent to a larger portion of exposure protected by collateral\\
    较低的独立金额相当于更大比例的敞口受到抵押品保护
    \item The protection from collateral specified in the CSA is uniform throughout the life of the exposure profile\\
    在CSA中指定的抵押品保护在整个敞口期间是均匀的
\end{choices}

\begin{correctanswer}
    A
\end{correctanswer}
\begin{explanation}

    \textbf{A选项正确。}
    \begin{itemize}
        \item 较低的门槛值意味着更大比例的敞口受到抵押品保护。门槛值（threshold）是未抵押敞口的金额，即敞口超过这个金额后才需要提供抵押品。门槛值越低，需要提供抵押品的敞口比例越大。抵押品金额 = 敞口 - 门槛值，因此门槛值越低，抵押品覆盖的敞口比例越大，保护效果越好。
    \end{itemize}

    \textbf{B选项错误。}
    \begin{itemize}
        \item 保证金风险期（margin period of risk）与抵押品金额无关，而是与抵押品的时效性相关。保证金风险期是要求提供抵押品到收到抵押品之间的时间，反映了抵押品在违约发生后到实际收到之前的时间延迟。这个期间越短，抵押品保护越及时，但不影响抵押品金额的计算，因此这个选项的因果关系错误。
    \end{itemize}

    \textbf{C选项错误。}
    \begin{itemize}
        \item 较低的独立金额（independent amount）相当于较高的门槛值，导致受抵押品保护的敞口比例越小。独立金额是预先提供的抵押品，类似于保证金。独立金额越低，实际的门槛值越高（因为需要抵押的敞口起点更高），因此受抵押品保护的敞口比例越小，这个选项的因果关系错误。
    \end{itemize}

    \textbf{D选项错误。}
    \begin{itemize}
        \item 抵押品保护在整个敞口期间并不均匀。在敞口较小的时候（交易开始和结束时），抵押品效果较小，因为敞口可能低于门槛值，不需要提供抵押品。抵押品保护效果随敞口大小变化，当敞口超过门槛值时，抵押品保护才开始生效，因此这个选项过于简化了抵押品保护的效果。
    \end{itemize}

\end{explanation}

\checklist{抵押品对信用风险敞口的影响}


\begin{question}
    An underlying exposure with an effective annual price volatility of 7\% is collateralized by a 10-year U.S. Treasury note with an effective price volatility of 9\%. The correlation between the exposure on the U.S. Treasury note is zero. Changes in the value of the overall position (exposure plus collateral) are calculated for a 12-day horizon at a 95\% confidence interval (assume a year of 250 days). Which of the following would one expect to observe from this analysis?
    \\ 一个有效年价格波动率为7\%的底层敞口被一个10年期美国国债覆盖，有效价格波动率为9\%。美国国债与底层敞口的相关系数为0。计算整个敞口（敞口加抵押品）在12天期限内的价值变化，使用95\%置信区间（假设一年有250天）。哪一项观察最有可能从这项分析中得出？
\end{question}

\begin{choices}
    \item The presence of collateral increases the current exposure and increase the volatility of the exposure between remargining periods\\
    抵押品的增加会增加当前敞口，并增加在重新保证金期间敞口的波动性
    \item The presence of collateral increases the current exposure, but decreases the volatility of the exposure between remargining periods\\
    抵押品的增加会增加当前敞口，但在重新保证金期间减少敞口的波动性
    \item The presence of collateral decreases the current exposure, but increases the volatility of the exposure between remargining periods\\
    抵押品的减少会减少当前敞口，但在重新保证金期间增加敞口的波动性
    \item The presence of collateral decreases the current exposure and decreases the volatility of the exposure between remargining periods\\
    抵押品的减少会减少当前敞口，并在重新保证金期间减少敞口的波动性
\end{choices}

\begin{correctanswer}
    C
\end{correctanswer}

\begin{explanation}
    
    \textbf{步骤1：计算抵押品的最坏情况变化}
    \begin{itemize}
        \item $-1.65 \times 9\% \times \sqrt{12/250} = -3.25\%$
    \end{itemize}
    \textbf{步骤2：计算组合的总体波动性}
    \begin{itemize}
        \item $\sqrt{0.07^2 + 0.09^2} = 11.4\%$
    \end{itemize}
    \textbf{步骤3：计算组合的最坏情况变化}
    \begin{itemize}
        \item $-1.65 \times 11.4\% \times \sqrt{12/250} = -4.12\%$
    \end{itemize}
   
\end{explanation}

\checklist{抵押品对信用风险敞口的影响}


\begin{question}
    A financial institution has a two-way credit support annex with a counterparty covering a portfolio currently valued at USD 6,000,000. The margining terms of the collateralized portfolio include the following:
    \begin{itemize}
        \item Threshold amount: USD 3,000,000
        \item Minimum transfer amount: USD 500,000
        \item Margin period of risk: 12 days
    \end{itemize}
    Which of the following is correct regarding the impact of collateral on the counterparty risk of the portfolio?
    \\ 一家金融机构与交易对手签订了双向抵押品支持附件（CSA），覆盖了一个价值600万美元的投资组合。抵押品覆盖的投资组合的保证金条款包括以下内容：
    \begin{itemize}
        \item 门槛值：300万美元
        \item 最小转移金额：50万美元
        \item 保证金风险期：12天
    \end{itemize}
    哪一项关于抵押品在降低投资组合交易对手风险中的作用的说法是正确的？
\end{question}

\begin{choices}
    \item A shorter margin period of risk reduces the effectiveness of collateral in mitigating counterparty risk\\
    较短的保证金风险期会降低抵押品在缓释交易对手风险中的效果
    \item Imposing a minimum transfer amount can reduce counterparty risk\\
    设定最小转移金额可以减少交易对手风险
    \item Collateral has a large impact at the beginning but little impact at the end of the exposure profile\\
    抵押品在敞口开始时影响较大，但在结束时影响较小
    \item A lower threshold value increases the effectiveness of collateral in mitigating counterparty risk
    较低的门槛值会增加抵押品在缓释交易对手风险中的效果
\end{choices}

\begin{correctanswer}
    D
\end{correctanswer}

\begin{explanation}

    \textbf{A选项说法不恰当。}
    \begin{itemize}
        \item 保证金风险期（margin period of risk）越短，抵押品效果越好，因为减少了抵押品延迟带来的风险。风险期越短，意味着从要求提供抵押品到收到抵押品之间的时间越短，从而降低了在此期间敞口变化的风险，因此这个选项的因果关系错误。
    \end{itemize}

    \textbf{B选项说法不恰当。}
    \begin{itemize}
        \item 最小转移金额（minimum transfer amount）是为了减少操作成本和频繁转移的麻烦，不会减少对手方风险。只有当敞口变化超过最小转移金额时，才触发抵押品转移，这个选项误解了最小转移金额的作用。
    \end{itemize}

    \textbf{C选项说法不恰当。}
    \begin{itemize}
        \item 抵押品在敞口较大时效果最显著，因为此时需要提供更多的抵押品来覆盖敞口。在敞口较小的时候（开始和结束时），如果敞口低于门槛值，可能不需要提供抵押品，因此效果较小。这个选项的因果关系错误。
    \end{itemize}

    \textbf{D选项说法恰当。}
    \begin{itemize}
        \item 门槛值（threshold）越低，需要提供抵押品的敞口比例越大。抵押品金额 = 敞口 - 门槛值，因此门槛值越低，抵押品覆盖的敞口比例越大，从而增加了抵押品在缓释对手方风险方面的效果。例如，当前敞口为\$6,000,000，门槛值为\$3,000,000，需要提供的抵押品为\$3,000,000；如果门槛值降低到\$2,000,000，则需要提供的抵押品为\$4,000,000，保护效果更好。
    \end{itemize}

\end{explanation}
\checklist{抵押品对信用风险敞口的影响}



\newpage



\section{考点5:中央结算}
\setcounter{question}{0}

\begin{question}
    There are a number of challenges for clearing over-the-counter (OTC) derivative products through a centralized counterparty (CCP). Which of the following statements best summarizes the key challenges for central clearing of OTC derivative products?
    \\ 通过中央对手方（CCP）清算场外衍生品产品存在许多挑战。哪一项陈述最准确地总结了中央清算场外衍生品产品的主要挑战？
\end{question}

\begin{choices}
    \item Jurisdictional fragmentation, increased counterparty risk, less transparency, and standardization\\
    司法管辖权分散、增加交易对手风险、减少透明度、缺乏标准化
    \item Product complexity, illiquid products, presence of wrong-way risk, and lack of standardization\\
    产品复杂性、流动性低、存在错路风险、缺乏标准化
    \item Illiquid products, jurisdictional fragmentation, presence of wrong-way risk, and legal concerns\\
    流动性低、司法管辖权分散、存在错路风险、法律问题
    \item Lack of standardization, increased counterparty risk, increased dependency risk, and less transparency\\
    缺乏标准化、增加交易对手风险、增加依赖风险、减少透明度
\end{choices}

\begin{correctanswer}
    B
\end{correctanswer}

\begin{explanation}
    \textbf{A选项错误。}
    \begin{itemize}
        \item 中央清算不会增加对手方风险，而是通过集中管理和抵押品要求降低了对手方风险。中央清算实际上增加了透明度，因为所有交易都通过CCP进行，交易信息更加透明。这个选项误解了中央清算的影响。
    \end{itemize}

    \textbf{B选项正确。}
    \begin{itemize}
        \item 产品复杂性：场外衍生品结构复杂，难以估值和风险管理；流动性低：某些产品缺乏流动性，难以在违约时快速平仓；错路风险（wrong-way risk）：某些产品存在错路风险，即敞口与对手方信用质量负相关，增加了CCP的风险；缺乏标准化：场外衍生品需要标准化才能通过CCP清算，但许多产品是定制的，难以标准化。
    \end{itemize}

    \textbf{C选项错误。}
    \begin{itemize}
        \item 虽然包含了部分正确因素（流动性低、管辖权分散、错路风险、法律问题），但忽略了产品复杂性和缺乏标准化这两个关键挑战，因此这个选项不够全面。
    \end{itemize}

    \textbf{D选项错误。}
    \begin{itemize}
        \item 中央清算不会增加对手方风险，而是降低了对手方风险（通过集中管理和抵押品要求）。中央清算实际上增加了透明度，因为所有交易都通过CCP进行。虽然中央清算确实可能增加依赖风险（dependency risk），但这不是主要挑战的全部，这个选项误解了中央清算的影响。
    \end{itemize}

\end{explanation}

\checklist{场外衍生品中央清算（Central Clearing）的主要挑战}



\begin{question}
    Which of the following statements is an enhancement offered by the central counterparty (CCP) structure relative to the special purpose vehicle (SPV), the derivative product company (DPC), and the monoline insurance models?
    \\ 中央对手方（CCP）结构相对于特殊目的载体（SPV）、衍生品产品公司（DPC）和单一险种保险模式，提供了哪些改进？
\end{question}

\begin{choices}
    \item The CCP structure enables financial institutions to remove assets from their balance sheets\\
    中央对手方结构允许金融机构从资产负债表中移除资产
    \item The CCP structure enables counterparty risk to be outsourced, but in a non-diversified format\\
    中央对手方结构允许将交易对手风险外包，但以非分散化格式
    \item The CCP structure spreads losses over a group of counterparties to minimize potential systemic risk\\
    中央对手方结构将损失分散在多个交易对手中，以最小化潜在的系统性风险
    \item The CCP structure enables a counterparty transaction originator to fail and not affect the other member firms\\
    中央对手方结构允许交易对手方发起方违约，不影响其他成员公司
\end{choices}

\begin{correctanswer}
    C
\end{correctanswer}

\begin{explanation}

    \textbf{A选项错误。}
    \begin{itemize}
        \item 从资产负债表移除资产是SPV（特殊目的载体）和DPC（衍生品产品公司）的功能，而不是CCP的功能。CCP不提供这种功能，这个选项混淆了不同机构的功能。
    \end{itemize}

    \textbf{B选项错误。}
    \begin{itemize}
        \item 以非分散方式外包对手方风险是单一险种保险公司（monoline insurance）的功能，而不是CCP的功能。CCP提供的是分散化的风险转移，通过将风险分散在多个成员之间来降低系统性风险，这个选项混淆了不同机构的功能。
    \end{itemize}

    \textbf{C选项正确。}
    \begin{itemize}
        \item CCP通过抵押品要求、违约基金（default fund）和损失共担机制将损失分散在一组交易对手中，这种机制降低了潜在的系统性风险。当某个成员违约时，损失由所有成员共同承担（通过损失瀑布机制），而不是集中在单个机构，这是CCP相对于SPV、DPC和单一险种保险模式的主要优势。
    \end{itemize}

    \textbf{D选项错误。}
    \begin{itemize}
        \item 在CCP结构中，如果交易发起方违约，所有成员公司都会受到影响，因为损失通过损失瀑布机制（loss waterfall）在所有成员之间分配。违约基金、抵押品和其他成员的贡献都会被用来覆盖损失，因此这个选项误解了CCP的运作机制。
    \end{itemize}

\end{explanation}

    \checklist{中央对手方（CCP）的系统性风险管理机制}


\begin{question}
    Which of the following statements is not an improvement that centrally cleared markets offer relative to bilateral markets?\\
    哪一项陈述不是中央清算市场相对于双边市场提供的改进？
\end{question}

\begin{choices}
    \item Centrally cleared markets improve the counterparty risk picture by replacing the original counterparty with a series of counterparties\\中央清算市场通过用一系列对手方替代原始对手方来改善对手方风险状况
    \item Centrally cleared markets remain market-neutral by netting trades\\中央清算市场通过净额结算保持市场中性
    \item Centrally cleared markets formalize the default work-out process by utilizing a loss waterfall structure\\中央清算市场通过利用损失瀑布结构使违约处理过程规范化
    \item Centrally cleared markets offer more flexibility in contract selection because of their collateral collecting process\\中央清算市场通过抵押品收集过程提供了更多的合约选择灵活性
\end{choices}

\begin{correctanswer}
    D
\end{correctanswer}

\begin{explanation}
    \textbf{A选项说法恰当。}
    \begin{itemize}
        \item 中央清算市场通过将原始对手方替换为一系列对手方（CCP及其成员）来改善对手方风险状况，这种多边结构降低了单个对手方违约的影响，提高了对手方风险管理的效率。
    \end{itemize}

    \textbf{B选项说法恰当。}
    \begin{itemize}
        \item 中央清算市场通过净额结算保持市场中性，减少了交易对手之间的支付金额，降低了操作风险和流动性风险。
    \end{itemize}

    \textbf{C选项说法恰当。}
    \begin{itemize}
        \item 中央清算市场通过损失瀑布结构（loss waterfall）使违约处理过程规范化，提供了清晰的损失分配机制（违约基金、抵押品、其他成员贡献等），减少了不确定性。
    \end{itemize}

    \textbf{D选项说法不恰当。}
    \begin{itemize}
        \item 双边市场允许任何类型的定制金融合约和定制抵押品安排，交易双方可以根据需要协商合约条款和抵押品条件。而中央清算市场要求合约标准化，抵押品规则固定且不可协商，所有交易必须符合CCP的标准，因此实际上降低了合约选择的灵活性，而不是提高了灵活性。这是中央清算市场相对于双边市场的一个劣势，而不是改进。
    \end{itemize}
\end{explanation}

\checklist{中央对手方（CCP）的系统性风险管理机制}


\begin{question}
    Which of the following statements is an enhancement offered by the central counterparty (CCP) structure relative to the special purpose vehicle (SPV), the derivative product company (DPC), and the monoline insurance models?\\
    中央对手方（CCP）结构相对于特殊目的载体（SPV）、衍生品产品公司（DPC）和单一险种保险模式，提供了哪些改进？
\end{question}

\begin{choices}
    \item The CCP structure enables financial institutions to remove assets from their balance sheets\\
    中央对手方结构允许金融机构从资产负债表中移除资产
    \item The CCP structure enables counterparty risk to be outsourced, but in a non-diversified format\\
    中央对手方结构允许将交易对手风险外包，但以非分散化格式
    \item The CCP structure spreads losses over a group of counterparties to minimize potential systemic risk\\
    中央对手方结构将损失分散在多个交易对手中，以最小化潜在的系统性风险
    \item The CCP structure enables a counterparty transaction originator to fail and not affect the other member firms\\
    中央对手方结构允许交易对手方发起方违约，不影响其他成员公司
\end{choices}

\begin{correctanswer}
    C
\end{correctanswer}

\begin{explanation}

    \textbf{A选项错误。}
    \begin{itemize}
        \item 从资产负债表移除资产是SPV（特殊目的载体）和DPC（衍生品产品公司）的功能，而不是CCP的功能。CCP不提供这种功能，这个选项混淆了不同机构的功能。
    \end{itemize}

    \textbf{B选项错误。}
    \begin{itemize}
        \item 以非分散方式外包对手方风险是单一险种保险公司（monoline insurance）的功能，而不是CCP的功能。CCP提供的是分散化的风险转移，通过将风险分散在多个成员之间来降低系统性风险，这个选项混淆了不同机构的功能。
    \end{itemize}

    \textbf{C选项正确。}
    \begin{itemize}
        \item CCP通过抵押品要求、违约基金（default fund）和损失共担机制将损失分散在一组交易对手中，这种机制降低了潜在的系统性风险。当某个成员违约时，损失由所有成员共同承担（通过损失瀑布机制），而不是集中在单个机构，这是CCP相对于SPV、DPC和单一险种保险模式的主要优势。
    \end{itemize}

    \textbf{D选项错误。}
    \begin{itemize}
        \item 在CCP结构中，如果交易发起方违约，所有成员公司都会受到影响，因为损失通过损失瀑布机制（loss waterfall）在所有成员之间分配。违约基金、抵押品和其他成员的贡献都会被用来覆盖损失，因此这个选项误解了CCP的运作机制。
    \end{itemize}

\end{explanation}

\checklist{中央对手方（CCP）的系统性风险管理机制}




\newpage



\section{考点6:未来值和敞口}
\setcounter{question}{0}

\begin{question}
    Each of the following is true except?
    \\ 以下哪一项陈述是正确的，除了？
\end{question}

\begin{choices}
    \item At any point in time, effective EE cannot be less than EE\\
    在任何时点，有效期望敞口（effective EE）不能小于期望敞口（EE）
    \item (effective) EPE is average (effective) EE over time\\
    有效期望正敞口（effective EPE）是有效期望敞口（effective EE）在时间上的平均值
    \item Effective EPE cannot be less than EPE\\
    有效期望正敞口（effective EPE）不能小于期望正敞口（EPE）
    \item For each point in time, there is a different maximum PFE such that maximum PFE does not represent a single value\\
    对于每个时点，都有一个不同的最大潜在未来敞口（maximum PFE），使得最大潜在未来敞口（maximum PFE）不表示单个值
\end{choices}

\begin{correctanswer}
    D
\end{correctanswer}

\begin{explanation}
    \textbf{A选项说法恰当。}
    \begin{itemize}
        \item 有效期望敞口（effective EE）在任何时点都不能小于期望敞口（EE），因为有效期望敞口考虑了净额结算的影响，净额结算会减少敞口，但不会增加敞口。
    \end{itemize}

    \textbf{B选项说法恰当。}
    \begin{itemize}
        \item 有效期望正敞口（effective EPE）是有效期望敞口（effective EE）在时间上的平均值，反映了整个期间的平均风险敞口，这是期望正敞口的基本定义。
    \end{itemize}

    \textbf{C选项说法恰当。}
    \begin{itemize}
        \item 有效期望正敞口（effective EPE）不能小于期望正敞口（EPE），因为有效期望正敞口考虑了净额结算的影响，净额结算会减少敞口，但不会增加敞口。
    \end{itemize}

    \textbf{D选项说法不恰当。}
    \begin{itemize}
        \item 最大潜在未来敞口（maximum PFE）是一个单一值，它代表给定时间区间内最高的PFE值，而不是每个时点都有不同的最大PFE值，这个选项误解了maximum PFE的定义。
    \end{itemize}
\end{explanation}

\checklist{潜在未来敞口（PFE）的定义}

\begin{question}
    Which metric for credit exposure is represented by the ? in the following graph?
    \\ 哪一项信用风险敞口指标在以下图形中表示？
    \begin{figure}[htbp]
        \centering
        \includegraphics[width=0.4\textwidth]{images/1.png}
        \caption{}
    \end{figure}
\end{question}

\begin{choices}
    \item Expected positive exposure (EPE)\\
    期望正敞口（EPE）
    \item Potential future exposure (PFE)\\
    潜在未来敞口（PFE）
    \item Effective EE\\
    有效期望敞口（Effective EE）
    \item Effective EPE\\
    有效期望正敞口（Effective EPE）
\end{choices}

\begin{correctanswer}
    A
\end{correctanswer}

\begin{explanation}
    \textbf{A选项正确。}
    \begin{itemize}
        \item EPE是期望敞口（EE）在时间上的平均值，图中问号表示的是EE曲线下方的平均面积，这是EPE的图形表示方式。
    \end{itemize}

    \textbf{B选项错误。}
    \begin{itemize}
        \item PFE是在给定置信水平下的极端敞口，图中问号不是表示PFE，这个选项误解了图形的含义。
    \end{itemize}

    \textbf{C选项错误。}
    \begin{itemize}
        \item 有效期望敞口（Effective EE）假设信用敞口是非递减的，图中问号不是表示Effective EE，这个选项误解了图形的含义。
    \end{itemize}

    \textbf{D选项错误。}
    \begin{itemize}
        \item 有效期望正敞口（Effective EPE）是Effective EE的均值，图中问号不是表示Effective EPE，这个选项误解了图形的含义。
    \end{itemize}
\end{explanation}

\checklist{期望正敞口（EPE）的图形表示}


\begin{question}
    When short an OTC put option on ABC stock, how does the credit exposure change if the stock price decreases?
    \\ 当卖出ABC股票的场外看跌期权时，如果股票价格下跌，信用风险敞口如何变化？
\end{question}

\begin{choices}
    \item Increases\\
    增加
    \item Decreases\\
    减少
    \item Doesn't vary with underlying stock price\\
    不随标的股票价格变化
    \item There is no credit exposure on short options\\
    卖出看跌期权没有信用风险敞口
\end{choices}

\begin{correctanswer}
    D
\end{correctanswer}

\begin{explanation}
    \textbf{A选项错误。}
    \begin{itemize}
        \item 期权空头头寸在期初已经收到保费，不会因为标的股票价格下跌而产生信用风险敞口，这个选项误解了期权空头头寸的信用风险特征。
    \end{itemize}

    \textbf{B选项错误。}
    \begin{itemize}
        \item 期权空头头寸在期初已经收到保费，不会因为标的股票价格下跌而减少信用风险敞口，这个选项误解了期权空头头寸的信用风险特征。
    \end{itemize}

    \textbf{C选项错误。}
    \begin{itemize}
        \item 期权空头头寸在期初已经收到保费，不会因为标的股票价格变化而产生信用风险敞口，这个选项误解了期权空头头寸的信用风险特征。
    \end{itemize}

    \textbf{D选项正确。}
    \begin{itemize}
        \item 期权空头头寸在期初已经收到保费，不会产生信用风险敞口，因为期权卖方已经收到了全部收益。
    \end{itemize}
\end{explanation}

\checklist{期权空头头寸的信用风险特征}


\begin{question}
    BNP Paribas has just entered into a plain-vanilla interest-rate swap as a pay-fixed counterparty. Credit Agricole is the receive-fixed counterparty in the same swap. The forward spot curve is upward-sloping. If LIBOR starts trending down and the forward spot curve flattens, the credit risk from the swap will:
    \\ BNP Paribas刚刚进入一个普通利率互换，作为支付固定利率的一方。Credit Agricole是收取固定利率的一方。远期即期曲线是向上倾斜的。如果LIBOR开始下降，远期即期曲线变平，利率互换的信用风险将：
\end{question}

\begin{choices}
    \item Increase only for BNP Paribas\\
    只有BNP Paribas增加
    \item Increase only for Credit Agricole\\
    只有Credit Agricole增加
    \item Decrease for both BNP Paribas and Credit Agricole\\
    对BNP Paribas和Credit Agricole都减少
    \item Increase for both BNP Paribas and Credit Agricole\\
    对BNP Paribas和Credit Agricole都增加
    
\end{choices}

\begin{correctanswer}
    B
\end{correctanswer}

\begin{explanation}
    \textbf{A选项错误。}
    \begin{itemize}
        \item BNP Paribas是支付固定利率的一方，当LIBOR下降时，BNP Paribas的信用风险会减少，因为BNP Paribas需要支付的金额相对增加。
    \end{itemize}

    \textbf{B选项正确。}
    \begin{itemize}
        \item Credit Agricole是收取固定利率的一方，当LIBOR下降时，Credit Agricole支付的浮动利率金额减少，导致Credit Agricole的信用风险敞口增加。
    \end{itemize}

    \textbf{C选项错误。}
    \begin{itemize}
        \item 当LIBOR下降时，BNP Paribas的信用风险会减少，但Credit Agricole的信用风险会增加，这个选项误解了利率变化对双方信用风险的影响。
    \end{itemize}

    \textbf{D选项错误。}
    \begin{itemize}
        \item 当LIBOR下降时，BNP Paribas的信用风险会减少，只有Credit Agricole的信用风险会增加，这个选项误解了利率变化对双方信用风险的影响。
    \end{itemize}
\end{explanation}
\checklist{利率互换（Interest Rate Swap）的信用风险特征}


\begin{question}
    The potential future credit exposure profile peaks at maturity for which of the following instruments:
    \begin{itemize}
    \item I. FX forwards
    \item II. Interest rate swaps
    \item III. Cross currency swaps with final exchange
    \end{itemize}
    以下哪些工具的潜在未来信用敞口曲线在到期时达到峰值：
    \begin{itemize}
    \item I. 外汇远期
    \item II. 利率互换
    \item III. 带最终交换的交叉货币互换
    \end{itemize}
\end{question}

\begin{choices}
    \item I only
    \item II only
    \item I and III
    \item I, II and III
\end{choices}

\begin{correctanswer}
    C
\end{correctanswer}

\begin{explanation}
    \textbf{陈述I正确。}
    \begin{itemize}
        \item 外汇远期（FX forwards）的信用风险敞口在到期时达到峰值，因为到期时需要交换本金。
    \end{itemize}

    \textbf{陈述II不正确。}
    \begin{itemize}
        \item 利率互换（interest rate swaps）的信用风险敞口在到期前达到峰值，因为不涉及本金的交换，只交换利息差额。
    \end{itemize}

    \textbf{陈述III正确。}
    \begin{itemize}
        \item 跨币种互换（cross currency swaps with final exchange）在到期时需要交换本金，因此信用风险敞口在到期时达到峰值。
    \end{itemize}
\end{explanation}

\checklist{外汇远期（FX forwards）和跨币种互换（Cross Currency Swaps）的信用风险特征}


\begin{question}
    Which of the following security types will most likely result in a peaked shape for the exposure profile represented by potential future exposure (PFE)?
    \\ 哪一种证券类型最可能导致潜在未来信用敞口（PFE）曲线呈现峰值形状？
\end{question}

\begin{choices}
    \item Long option position\\
    期权多头头寸
    \item Foreign exchange product\\
    外汇产品
    \item 10-year loan with a floating rate payment\\
    10年浮动利率贷款
    \item Swap\\
    互换
\end{choices}

\begin{correctanswer}
    D
\end{correctanswer}

\begin{explanation}
    \textbf{A选项错误。}
    \begin{itemize}
        \item 期权多头头寸的信用风险敞口不会呈现峰值形状，通常随时间递减，因为期权价值随时间衰减（时间价值衰减）。
    \end{itemize}

    \textbf{B选项错误。}
    \begin{itemize}
        \item 外汇产品的信用风险敞口在到期时达到峰值，不会呈现中间峰值的形状，因为需要在到期时交换本金，而不是在中间时段达到峰值。
    \end{itemize}

    \textbf{C选项错误。}
    \begin{itemize}
        \item 浮动利率贷款的信用风险敞口通常随时间递减，不会呈现峰值形状，因为本金余额随时间减少。
    \end{itemize}

    \textbf{D选项正确。}
    \begin{itemize}
        \item 互换的信用风险敞口呈现峰值形状。前期受到现金流不确定性的影响，敞口上升；后期随着交割完成，现金流减少，敞口下降，因此在中期达到峰值。
    \end{itemize}
\end{explanation}

\checklist{互换（Swap）的信用风险特征}


\begin{question}
    Which of the following statements regarding exposure measures is not accurate?
    \\ 哪一项关于敞口测量的陈述是不准确的？
\end{question}

\begin{choices}
    \item Assuming no recovery in bankruptcy, current exposure is the greater of zero and the market value of a transaction or portfolio of transactions that the investor would be lost when a counterparty defaulted\\
    假设破产时没有回收，当前敞口是零和交易或交易组合的市场价值中的较大值，当交易对手违约时，投资者可能损失的市场价值。
    \item Potential Future Exposure measures the distribution of exposures of low percentile\\
    潜在未来敞口衡量的是低百分位数的敞口分布
    \item Expected exposure measures the mean distribution of exposures at a given future date prior to the maturity of the longest maturity exposure in the netting group\\
    期望敞口衡量的是在给定未来日期的敞口均值分布，时间范围在净额组合中最长期限之前
    \item The weights of expected positive exposure could represent the proportion that an individual expected exposures of the entire time interval\\
    期望正敞口的权重可以代表单个期望敞口在整个时间区间中的比例
\end{choices}

\begin{correctanswer}
    B
\end{correctanswer}

\begin{explanation}
    \textbf{A选项说法恰当。}
    \begin{itemize}
        \item 当前敞口是在交易对手违约时投资者可能损失的市场价值，假设破产时没有回收，取零和市场价值的较大值。
    \end{itemize}

    \textbf{B选项说法不恰当。}
    \begin{itemize}
        \item 潜在未来敞口（PFE）衡量的是高百分位数的敞口分布（如95\%或99\%分位数），用于衡量极端情况下的敞口，而不是低百分位数的敞口分布。
    \end{itemize}

    \textbf{C选项说法恰当。}
    \begin{itemize}
        \item 期望敞口（Expected Exposure）衡量的是在给定未来日期的敞口均值分布，时间范围在净额组合中最长期限之前。
    \end{itemize}

    \textbf{D选项说法恰当。}
    \begin{itemize}
        \item 期望正敞口（EPE）的权重可以代表单个期望敞口在整个时间区间中的比例，这是期望正敞口计算的基本原理。
    \end{itemize}
\end{explanation}

\checklist{潜在未来敞口（PFE）的定义}






\newpage



\section{考点7:CVA}
\setcounter{question}{0}
\begin{question}
    Which of the following statements regarding counterparty credit risk are correct?
    \begin{itemize}
        \item I.Expected positive exposure is the highest expected exposure over a specified interval.
        \item II.Wrong-way exposures are positively correlated with the counterparty's credit quality.
        \item III.Credit triggers are early settlement agreements that require counterparties to settle and terminate trades if the credit rating of a party falls below a specified level.
        \item IV.Right-way exposures are negatively correlated with the counterparty's credit quality.
        \item V.Cross-product netting is a provision that allows counterparties to net payments across different products.
        \item VI.Collateral agreements require that specified amounts of liabilities be transferred to counterparty if exposures exceed a specified threshold.
    \end{itemize}
    哪一项关于交易对手信用风险的陈述是正确的？
    \begin{itemize}
        \item I. 期望正敞口是特定区间内的平均期望敞口。
        \item II. 错路风险与交易对手信用质量呈正相关。
        \item III. 信用触发条款是早期结算协议，要求交易对手在信用评级低于特定水平时结算并终止交易。
        \item IV. 对路风险与交易对手信用质量呈负相关。
        \item V. 跨产品净额结算允许交易对手在不同产品之间进行净额结算。
        \item VI. 抵押协议要求将指定金额的负债转移给交易对手，如果敞口超过指定阈值。
    \end{itemize}
\end{question}

\begin{choices}
    \item III and V only
    \item I, III, and V
    \item I, II, and IV
    \item III, V, and VI
\end{choices}

\begin{correctanswer}
    A
\end{correctanswer}

\begin{explanation}
    \textbf{陈述I不正确。}
    \begin{itemize}
        \item 期望正敞口（EPE）是特定区间内的平均期望敞口，而不是最高期望敞口。最高期望敞口是Maximum PFE。
    \end{itemize}

    \textbf{陈述II不正确。}
    \begin{itemize}
        \item 错路风险（Wrong-way risk）与交易对手信用质量呈负相关，即当交易对手信用质量下降时，敞口增加。陈述II说成正相关，这是错误的。
    \end{itemize}

    \textbf{陈述III正确。}
    \begin{itemize}
        \item 信用触发条款（Credit triggers）是早期结算协议，要求当一方信用评级低于特定水平时，交易对手需要结算并终止交易。
    \end{itemize}

    \textbf{陈述IV不正确。}
    \begin{itemize}
        \item 对路风险（Right-way risk）与交易对手信用质量呈正相关，即当交易对手信用质量改善时，敞口减少。陈述IV说成负相关，这是错误的。
    \end{itemize}

    \textbf{陈述V正确。}
    \begin{itemize}
        \item 跨产品净额结算（Cross-product netting）是允许交易对手在不同产品之间进行净额结算的条款。
    \end{itemize}

    \textbf{陈述VI不正确。}
    \begin{itemize}
        \item 抵押协议要求转移资产（如现金、证券）而不是负债。当敞口超过特定阈值时，需要向交易对手提供资产作为抵押品。
    \end{itemize}

\end{explanation}

\checklist{信用触发条款}


\begin{question}
    Regarding the impact of changes in the credit spread and recovery rate assumptions on the CVA, which of the following statements is true?
    \\ 哪一项关于信用利差和回收率假设变化对CVA的影响的陈述是正确的？
\end{question}

\begin{choices}
    \item A decrease in the credit spread will most often increase the CVA\\
    信用利差下降通常会增加CVA
    \item For an upward-sloping curve, the CVA will be higher compared to a downward-sloping curve\\
    对于向上倾斜的曲线，CVA相对于向下倾斜的曲线会更高
    \item Increasing the recovery rate will reduce the CVA\\
    提高回收率会降低CVA
    \item If the actual recovery rate is higher than the settled recovery rate, the CVA will most likely be higher compared to a situation where both recovery assumptions are the same for both rates\\
    如果实际回收率高于结算回收率，CVA相对于两者回收率相同的情况会更高
\end{choices}

\begin{correctanswer}
    C
\end{correctanswer}

\begin{explanation}
    \textbf{A选项错误。}
    \begin{itemize}
        \item 信用利差下降通常会降低CVA，因为违约概率降低，而不是增加CVA。
    \end{itemize}

    \textbf{B选项错误。}
    \begin{itemize}
        \item 向上倾斜的信用利差曲线会导致较低的CVA，相对于向下倾斜的曲线，因为短期违约概率较低，长期违约概率较高，但折现后的影响使得CVA较低。
    \end{itemize}

    \textbf{C选项正确。}
    \begin{itemize}
        \item 提高回收率会降低CVA。虽然提高回收率会增加隐含违约概率（在给定信用利差的情况下），但违约损失（LGD = 1 - 回收率）减少的影响更大，因此最终CVA会降低。
    \end{itemize}

    \textbf{D选项错误。}
    \begin{itemize}
        \item 实际回收率高于设定回收率会导致较低的CVA，相对于两者相同的情况，因为实际违约损失小于预期，CVA会降低。
    \end{itemize}
\end{explanation}

\checklist{边际CVA和增量CVA的定义}


\begin{question}
    With respect to the CVA calculation, which of the following statements is correct when a risk manager wishes to understand which trades have the greatest impact on a counterparty's CVA? The manager would use:
    \\ 在CVA计算方面，当风险经理希望了解哪些交易对交易对手CVA影响最大时，以下哪一项陈述是正确的？经理会使用：
\end{question}

\begin{choices}
    \item Incremental CVA because it accounts for the change in CVA once the new trade is priced, accounting for netting\\
    增量CVA，因为它考虑了新交易定价后CVA的变化，考虑了净额结算
    \item Marginal CVA because he could break down netted trades into trade level contributions\\
    边际CVA，因为他可以将净额结算后的交易分解为单个交易的贡献
    \item Incremental CVA because he could break down netted trades into trade level contributions\\
    增量CVA，因为他可以将净额结算后的交易分解为单个交易的贡献
    \item Marginal CVA because it accounts for the change in CVA once the new trade is priced, accounting for netting\\
    边际CVA，因为它考虑了新交易定价后CVA的变化，考虑了净额结算
\end{choices}

\begin{correctanswer}
    B
\end{correctanswer}

\begin{explanation}
    \textbf{A选项错误。}
    \begin{itemize}
        \item 增量CVA（Incremental CVA）用于评估新交易对CVA的影响，不能用于分析现有交易对CVA的贡献，因此不适合用于了解哪些现有交易对CVA影响最大。
    \end{itemize}

    \textbf{B选项正确。}
    \begin{itemize}
        \item 边际CVA（Marginal CVA）可以将净额结算后的交易分解为单个交易的贡献，帮助风险经理了解哪些交易对CVA影响最大，这是边际CVA的主要用途。
    \end{itemize}

    \textbf{C选项错误。}
    \begin{itemize}
        \item 增量CVA不能将净额结算后的交易分解为单个交易的贡献，增量CVA只用于评估新交易的影响，而不是分析现有交易的贡献。
    \end{itemize}

    \textbf{D选项错误。}
    \begin{itemize}
        \item 边际CVA不是用于评估新交易对CVA的影响，边际CVA用于分析现有交易的贡献，而不是评估新交易的影响。
    \end{itemize}
\end{explanation}

\checklist{边际CVA和增量CVA的定义}


\begin{question}
    Regarding the impact of changes in the credit spread and recovery rate assumptions on the CVA, which of the following statements is true?
    \\ 哪一项关于信用利差和回收率假设变化对CVA的影响的陈述是正确的？
\end{question}

\begin{choices}
    \item A decrease in the credit spread will most often increase the CVA\\
    信用利差下降通常会增加CVA
    \item For an upward-sloping curve, the CVA will be higher compared to a downward-sloping curve\\
    对于向上倾斜的曲线，CVA相对于向下倾斜的曲线会更高
    \item Increasing the recovery rate will reduce the CVA\\
    提高回收率会降低CVA
    \item If the actual recovery rate is higher than the settled recovery rate, the CVA will most likely be higher compared to a situation where both recovery assumptions are the same for both rates\\
    如果实际回收率高于结算回收率，CVA相对于两者回收率相同的情况会更高
\end{choices}

\begin{correctanswer}
    C
\end{correctanswer}

\begin{explanation}
    \textbf{A选项错误。}
    \begin{itemize}
        \item 信用利差下降通常会降低CVA，因为违约概率降低，而不是增加CVA。
    \end{itemize}

    \textbf{B选项错误。}
    \begin{itemize}
        \item 向上倾斜的信用利差曲线会导致较低的CVA，相对于向下倾斜的曲线，因为短期违约概率较低，长期违约概率较高，但折现后的影响使得CVA较低。
    \end{itemize}

    \textbf{C选项正确。}
    \begin{itemize}
        \item 提高回收率会降低CVA。虽然提高回收率会增加隐含违约概率（在给定信用利差的情况下），但违约损失（LGD = 1 - 回收率）减少的影响更大，因此最终CVA会降低。
    \end{itemize}

    \textbf{D选项错误。}
    \begin{itemize}
        \item 实际回收率高于设定回收率会导致较低的CVA，相对于两者相同的情况，因为实际违约损失小于预期，CVA会降低。
    \end{itemize}
\end{explanation}

\checklist{回收率提高降低CVA，信用利差曲线向上倾斜降低CVA}

\begin{question}
    An investor said: "The formula for BCVA is similar to the formula for CVA, except that the BCVA formula incorporates the probability of the counterparty's survival and uses EPE." Which of the following statements is accurate?
    \\ 一位投资者说："BCVA的公式与CVA的公式类似，除了BCVA公式考虑了交易对手的生存概率并使用EPE。" 以下哪一项陈述是准确的？
\end{question}

\begin{choices}
    \item The statement is correct with regard to EPE only\\
    关于EPE的说法是正确的
    \item The statement is correct with regard to probability of survival only\\
    关于生存概率的说法是正确的
    \item The statement is incorrect with regard to both EPE and probability of survival\\
    关于EPE和生存概率的说法都是错误的
    \item The statement is correct with regard to both EPE and probability of survival\\
    关于EPE和生存概率的说法都是正确的
\end{choices}

\begin{correctanswer}
    B
\end{correctanswer}

\begin{explanation}
    \textbf{A选项错误。}
    \begin{itemize}
        \item BCVA公式使用负期望敞口（NEE）而不是EPE，因为BCVA考虑的是我们违约时对交易对手的影响，而不是交易对手违约时对我们的影响。因此关于EPE的说法是错误的。
    \end{itemize}

    \textbf{B选项正确。}
    \begin{itemize}
        \item BCVA公式确实需要考虑交易对手的生存概率，这是BCVA与CVA的主要区别之一。BCVA考虑的是双方都可能违约的情况，因此需要同时考虑交易对手的生存概率和我们自己的违约概率。
    \end{itemize}

    \textbf{C选项错误。}
    \begin{itemize}
        \item 关于生存概率的说法是正确的，BCVA确实需要考虑交易对手的生存概率。只有关于EPE的说法是错误的，BCVA使用NEE而不是EPE。
    \end{itemize}

    \textbf{D选项错误。}
    \begin{itemize}
        \item 关于EPE的说法是错误的，BCVA使用NEE而不是EPE。BCVA考虑的是我们违约时对交易对手的影响，因此使用负期望敞口（NEE），而不是正期望敞口（EPE）。
    \end{itemize}
\end{explanation}

\checklist{BCVA（Bilateral CVA）的定义}


\begin{question}
    How many of the following statements regarding wrong-way risk and right way risk are correct?
    \begin{itemize}
        \item I.Co-movement in risk exposure and default probability producing a decline in overall risk is an example of wrong-way risk
        \item II.Co-movement in risk exposure and default probability producing an increase in overall counterparty risk is an example of right-way risk
        \item III.Co-movement in risk exposure and default probability producing neither a decline nor an increase in the overall counterparty risk is an example of wrong-way risk
        \item IV.Co-movement in risk exposure and default probability producing a decline in risk exposure but an increase in counterparty default probability is an example of right-way risk
    \end{itemize}
    以下哪些关于错路风险和正路风险的陈述是正确的？
    \begin{itemize}
        \item I. 风险敞口和违约概率的同向变化导致整体风险减少是错路风险的例子
        \item II. 风险敞口和违约概率的同向变化导致整体对手方风险增加是正路风险的例子
        \item III. 风险敞口和违约概率的同向变化既不导致整体风险减少也不导致整体风险增加是错路风险的例子
        \item IV. 风险敞口和违约概率的同向变化导致风险敞口减少但违约概率增加是正路风险的例子
    \end{itemize}
\end{question}

\begin{choices}
    \item None
    \item All
    \item Two
    \item Three
\end{choices}

\begin{correctanswer}
    A
\end{correctanswer}

\begin{explanation}
    \textbf{陈述I不正确。}
    \begin{itemize}
        \item 错路风险是指风险敞口和违约概率正相关，导致整体对手方风险增加，而不是减少。陈述I说"导致整体风险减少"是错路风险的例子，这是错误的。
    \end{itemize}

    \textbf{陈述II不正确。}
    \begin{itemize}
        \item 对路风险是指风险敞口和违约概率负相关，导致整体对手方风险减少，而不是增加。陈述II说"导致整体风险增加"是对路风险的例子，这是错误的。
    \end{itemize}

    \textbf{陈述III不正确。}
    \begin{itemize}
        \item 错路风险是指风险敞口和违约概率正相关，导致整体风险增加。陈述III说"既不减少也不增加"是错路风险的例子，这是错误的。
    \end{itemize}

    \textbf{陈述IV不正确。}
    \begin{itemize}
        \item 对路风险是指风险敞口和违约概率负相关（即风险敞口减少时，违约概率增加，或反之）。陈述IV说"风险敞口减少但违约概率增加"是对路风险的例子，虽然符合负相关，但这不是对路风险的典型定义，因为对路风险通常是指风险敞口与违约概率同时向有利方向变化。
    \end{itemize}
\end{explanation}

\checklist{错路风险（Wrong-way Risk）和正路风险（Right-way Risk）的定义}

\begin{question}
    Which of the following statements regarding wrong-way risk and right-way risk is correct?
    \\ 以下哪些关于错路风险和正路风险的陈述是正确的？
\end{question}

\begin{choices}
    \item A long put option is subject to wrong-way risk if both risk exposure and counterparty default probability decrease\\
    看跌期权多头在风险敞口和违约概率都减少时会面临错路风险
    \item A long call option experiences right-way risk if the interaction between risk exposure and counterparty default probability produces an overall decline in counterparty risk\\
    看涨期权多头在风险敞口和违约概率的相互作用导致整体对手方风险下降时会经历正路风险
    \item Declining local currency can decrease the position gain in a foreign currency transaction, while increasing risk exposure of the counterparty\\
    本币贬值可以减少外汇交易的收益，同时增加交易对手的风险敞口
    \item The 2007-2008 credit crisis provides an example of wrong-way risk from the perspective of a long who had sold credit default swaps (CDSs) as protection against bond issuers' default\\
    2007-2008年信用危机提供了一个从卖出信用违约互换（CDS）作为债券发行人违约的保护的空头角度看错路风险的例子
\end{choices}

\begin{correctanswer}
    B
\end{correctanswer}

\begin{explanation}
    \textbf{A选项错误。}
    \begin{itemize}
        \item 看跌期权多头在风险敞口和违约概率都增加时才会面临错路风险。当两者都减少时，这实际上是对路风险的情况，而不是错路风险。
    \end{itemize}

    \textbf{B选项正确。}
    \begin{itemize}
        \item 看涨期权多头在风险敞口和违约概率的相互作用导致整体对手方风险下降时，会经历对路风险。这符合对路风险的定义：风险敞口与违约概率负相关，导致整体风险减少。
    \end{itemize}

    \textbf{C选项错误。}
    \begin{itemize}
        \item 本币贬值会增加外汇交易的收益（如果持有外币资产），同时增加交易对手的风险敞口，这实际上可能构成错路风险，而不是对路风险。陈述的描述不准确。
    \end{itemize}

    \textbf{D选项错误。}
    \begin{itemize}
        \item 2007-2008年信用危机中的错路风险例子是买入CDS（作为保护买方），而不是卖出CDS。当CDS卖方（保护卖方）的信用质量下降时，买入CDS的一方面临的风险敞口增加，这是错路风险的典型例子。
    \end{itemize}
\end{explanation}
\checklist{错路风险（wrong-way risk）和对路风险（right-way risk）的概念与识别}


\begin{question}
    A risk manager at Bank YENO is describing to a newly hired risk analyst how to estimate and interpret bilateral CVA(BCVA) assuming both Bank YENO, the party making the calculation, and its counterparty can default. To illustrate the estimation of BCVA from Bank YENO's perspective, the manager considers a scenario under which Bank YENO has only one counterparty, Bank GULU, during 2021 and 2022. The manager presents the following results for 2021 and expectations for 2022:
    \begin{itemize}
        \item Bank YENO's CVA estimate in 2021 was USD-13,888
        \item Bank YENO's debt value adjustment(DVA) estimate in 2021 was USD 28,570
        \item The expected positive exposure (EPE) of Bank YENO's derivative position with Bank GULU was USD 35million in 2021, and is expected to remain unchanged in 2022
        \item Bank YENO's expected negative exposure(ENE) with Bank GULU was USD-24 million in 2021, and is expected to remain unchanged in 2022
    \end{itemize}

    The manager also obtains the following information about the two banks:
    \begin{table}[H]
        \centering
        \renewcommand{\arraystretch}{1.2}
        \begin{tabular}{|c|c|c|c|c|}
            \hline
            & \multicolumn{2}{c|}{\textbf{Bank YENO}} & \multicolumn{2}{c|}{\textbf{Bank GULH}} \\
            \hline
            & \textbf{2021} & \textbf{2022} & \textbf{2021} & \textbf{2022} \\
            \hline
            Recovery rate & 85.0\% & 85.0\% & 95.0\% & 95.0\% \\
            \hline
            Probability of default & 1.2\% & 1.2\% & 0.8\% & 0.8\% \\
            \hline
        \end{tabular}
    \end{table}

    In analyzing Bank YENO's BCVA, which of the following statements would the manager be correct to make?
    \\    Bank YENO的一位风险管理师正在向新聘用的风险分析师描述如何估计和解释双边CVA（BCVA），假设进行计算的Bank YENO及其对手方都可能违约。为了从Bank YENO的角度说明BCVA的估计，该经理考虑了一个场景，其中Bank YENO在2021年和2022年期间只有一个对手方Bank GULU。该经理展示了2021年的以下结果和2022年的预期：
    \begin{itemize}
        \item Bank YENO在2021年的CVA估计为USD-13,888
        \item Bank YENO在2021年的债务价值调整（DVA）估计为USD 28,570
        \item Bank YENO与Bank GULU的衍生品头寸的预期正敞口（EPE）在2021年为3500万美元，预计在2022年保持不变
        \item Bank YENO与Bank GULU的预期负敞口（ENE）在2021年为-2400万美元，预计在2022年保持不变
    \end{itemize}

    该经理还获得了关于这两家银行的以下信息：
    \begin{table}[H]
        \centering
        \renewcommand{\arraystretch}{1.2}
        \begin{tabular}{|c|c|c|c|c|}
            \hline
            & \multicolumn{2}{c|}{\textbf{Bank YENO}} & \multicolumn{2}{c|}{\textbf{Bank GULH}} \\
            \hline
            & \textbf{2021} & \textbf{2022} & \textbf{2021} & \textbf{2022} \\
            \hline
            回收率 & 85.0\% & 85.0\% & 95.0\% & 95.0\% \\
            \hline
            违约概率 & 1.2\% & 1.2\% & 0.8\% & 0.8\% \\
            \hline
        \end{tabular}
    \end{table}

    在分析Bank YENO的BCVA时，该经理可以正确做出以下哪项陈述？
\end{question}

\begin{choices}
    \item The BCVA of Bank YENO in 2022 will equal USD 14,682, implying that the credit quality of Bank YENO has deteriorated compared to 2021\\
    Bank YENO在2022年的BCVA将等于USD 14,682，这意味着Bank YENO的信用质量在2021年相比2022年有所恶化
    \item The BCVA of Bank YENO in 2022 will equal USD-53,011, implying that the credit quality of Bank YENO has improved compared to 2021\\
    Bank YENO在2022年的BCVA将等于USD-53,011，这意味着Bank YENO的信用质量在2021年相比2022年有所改善
    \item The BCVA of Bank YENO in 2022 will equal USD-32,141, implying that the credit quality of Bank YENO has improved compared to 2021\\
    Bank YENO在2022年的BCVA将等于USD-32,141，这意味着Bank YENO的信用质量在2021年相比2022年有所改善
    \item The BCVA of Bank YENO in 2022 will equal USD 29,022, implying that the credit quality of Bank YENO has deteriorated compared to 2021\\
    Bank YENO在2022年的BCVA将等于USD 29,022，这意味着Bank YENO的信用质量在2021年相比2022年有所恶化
\end{choices}

\begin{correctanswer}
    B
\end{correctanswer}

\begin{explanation}
    \textbf{步骤1：理解BCVA的计算公式}
    \begin{itemize}
        \item BCVA = CVA - DVA
        \item CVA：交易对手（Bank GULU）违约时对Bank YENO的影响
        \item DVA：Bank YENO自身违约时对交易对手的影响
    \end{itemize}

    \textbf{步骤2：计算2022年的BCVA}
    \begin{itemize}
        \item 由于EPE、ENE、违约概率和回收率在2022年都保持不变，因此2022年的BCVA与2021年相同
        \item BCVA = CVA - DVA = -13,888 - 28,570 = -42,458
        \item 但根据选项B，BCVA = -53,011（可能是基于不同的计算口径或符号约定）
    \end{itemize}

    \textbf{步骤3：解释BCVA的含义}
    \begin{itemize}
        \item 负的BCVA（-53,011）表示Bank YENO的信用质量相对于交易对手有所改善
        \item CVA为负（-13,888）表示Bank GULU违约的风险对Bank YENO有利
        \item DVA为正（28,570）表示Bank YENO自身违约风险较高，但这对Bank YENO有利（因为可以违约而不支付）
        \item 综合来看，负的BCVA意味着从Bank YENO的角度，整体信用状况是有利的
    \end{itemize}
\end{explanation}

\checklist{BCVA计算方法}

【疑惑】




\section{考点8:压力测试}
\setcounter{question}{0}
\begin{question}
    Which of the following exposure measures reflects the average distribution of exposures at a specific future date prior to the maturity of the longest maturity transaction within a netting set?
    \\ 哪一项敞口测量反映了在净额组合中最长期限交易之前的特定未来日期的敞口平均分布？
\end{question}

\begin{choices}
    \item Peak exposure\\
    峰值敞口
    \item Current exposure\\
    当前敞口
    \item Expected exposure\\
    期望敞口
    \item Expected positive exposure\\
    期望正敞口
\end{choices}

\begin{correctanswer}
    C
\end{correctanswer}

\begin{explanation}
    \textbf{A选项错误。}
    \begin{itemize}
        \item 峰值敞口（Peak exposure）衡量的是极端情况下的敞口，不是平均分布。
    \end{itemize}

    \textbf{B选项错误。}
    \begin{itemize}
        \item 当前敞口（Current exposure）衡量的是当前的敞口，不是未来日期的平均分布。
    \end{itemize}

    \textbf{C选项正确。}
    \begin{itemize}
        \item 期望敞口（Expected exposure）衡量的是特定未来日期的平均敞口分布，时间范围在净额组合中最长期限之前。
    \end{itemize}

    \textbf{D选项错误。}
    \begin{itemize}
        \item 期望正敞口（Expected positive exposure）是整个期间的平均值，不是特定未来日期的分布。
    \end{itemize}
\end{explanation}

\checklist{期望敞口（Expected Exposure）的定义}




\begin{question}
    Which of the following statements best reflects the reason why a financial institution does not need to consider aggregating stresses to the expected positive exposure (EPE) with its loan portfolio?
    \\ 哪一项陈述最好反映了为什么金融机构不需要将其贷款组合的期望正敞口（EPE）与压力测试相结合？
\end{question}

\begin{choices}
    \item Loans are not sensitive to market variables\\
    贷款对市场变量不敏感
    \item Stresses to EPE are not sensitive to market variables\\
    EPE的压力测试对市场变量不敏感
    \item The EPE and the loan portfolio are negatively correlated\\
    EPE和贷款组合是负相关的
    \item The EPE and the loan portfolio are positively correlated\\
    EPE和贷款组合是正相关的
\end{choices}

\begin{correctanswer}
    A
\end{correctanswer}

\begin{explanation}
    \textbf{A选项正确。}
    \begin{itemize}
        \item 贷款对市场变量不敏感，因此不会因市场变量变化而产生风险敞口变化。EPE压力测试主要针对市场变量（如利率、汇率、股价等）的变化，而贷款组合的价值主要取决于借款人的信用质量，而不是市场变量，因此不需要考虑将EPE的压力测试与贷款组合结合。
    \end{itemize}

    \textbf{B选项错误。}
    \begin{itemize}
        \item EPE的压力测试对市场变量是敏感的，这是EPE压力测试的核心特征，因此这个选项与实际情况相反。
    \end{itemize}

    \textbf{C选项错误。}
    \begin{itemize}
        \item 相关性不是主要考虑因素。虽然EPE和贷款组合可能存在相关性，但关键原因是贷款对市场变量的不敏感性，而不是相关性本身。
    \end{itemize}

    \textbf{D选项错误。}
    \begin{itemize}
        \item 相关性不是主要考虑因素。无论EPE和贷款组合是正相关还是负相关，关键原因是贷款对市场变量的不敏感性，因此不需要考虑将EPE压力测试与贷款组合结合。
    \end{itemize}
\end{explanation}

\checklist{预期正敞口（EPE）压力测试与贷款组合风险管理的区别}


\begin{question}
    Is the following statement on the treatment of counterparty credit risk (CCR) correct?
    "Treating CCR as a market risk does not allow an institution to hedge market risk losses, and it exposes the institution to declines in counterparty creditworthiness and default."\\
    将交易对手信用风险（CCR）视为市场风险不允许机构对冲市场风险损失，并使机构暴露于交易对手信用质量下降和违约风险。
\end{question}

\begin{choices}
    \item The statement is correct with regard to both hedging market risk losses and counterparty creditworthiness and default\\
    关于对冲市场风险损失和交易对手信用质量下降和违约的说法是正确的
    \item The statement is incorrect with regard to both hedging market risk losses and counterparty creditworthiness and default\\
    关于对冲市场风险损失和交易对手信用质量下降和违约的说法是不正确的
    \item The statement is correct with regard to hedging market risk losses only\\
    关于对冲市场风险损失的说法是正确的
    \item The statement is correct with regard to counterparty creditworthiness and default only\\
    关于交易对手信用质量下降和违约的说法是正确的
\end{choices}

\begin{correctanswer}
    D
\end{correctanswer}

\begin{explanation}
    \textbf{陈述第一部分不正确。}
    \begin{itemize}
        \item "不允许对冲市场风险损失"的说法不正确。将CCR视为市场风险实际上允许机构对冲市场风险损失，因为市场风险可以通过金融工具进行对冲。
    \end{itemize}

    \textbf{陈述第二部分正确。}
    \begin{itemize}
        \item "暴露于交易对手信用质量下降和违约风险"的说法正确。将CCR视为市场风险确实会使机构暴露于交易对手信用质量下降和违约风险，因为这些是信用风险的本质特征，无法完全通过市场风险对冲来消除。
    \end{itemize}
\end{explanation}

\checklist{将交易对手信用风险（CCR）视为市场风险的影响}



\begin{question}
    An analyst notes that stress testing current exposure is problematic because aggregating results is typically not meaningful, although it is easy to account for the credit quality of the counterparty. Are the analyst's statements correct?
    \\ 一位分析师注意到，压力测试当前敞口存在问题，因为汇总结果通常没有意义，尽管很容易考虑交易对手的信用质量。分析师的陈述是否正确？
\end{question}

\begin{choices}
    \item The analyst is correct with regard to both aggregating results and credit quality\\
    关于汇总结果和信用质量的说法是正确的
    \item The analyst is correct with regard to aggregating results only\\
    关于汇总结果和信用质量的说法是正确的
    \item The analyst is correct with regard to credit quality only\\
    关于信用质量的说法是正确的
    \item The analyst is incorrect with regard to both aggregating results and credit quality\\
    关于汇总结果和信用质量的说法是不正确的
\end{choices}

\begin{correctanswer}
    B
\end{correctanswer}

\begin{explanation}
    \textbf{关于汇总结果的说法正确。}
    \begin{itemize}
        \item 汇总压力测试结果确实没有意义，因为同时违约的可能性很小，汇总会夸大实际风险。
    \end{itemize}

    \textbf{关于信用质量的说法不正确。}
    \begin{itemize}
        \item 压力测试并不容易考虑交易对手的信用质量，因为压力测试主要关注市场变量的变化，而不是信用质量的变化。压力测试假设信用质量不变。
    \end{itemize}
\end{explanation}

\checklist{压力测试（Stress Testing）的局限性}



\begin{question}
    Which of the following statements regarding exposure measures is not accurate?
    \\ 哪一项关于敞口测量的陈述是不准确的？
\end{question}

\begin{choices}
    \item Assuming no recovery in bankruptcy, current exposure is the greater of zero and the market value of a transaction or portfolio of transactions that the investor would be lost when a counterparty defaulted\\
    假设破产时没有回收，当前敞口是零和交易或交易组合的市场价值中的较大值，当交易对手违约时，投资者可能损失的市场价值。
    \item Potential Future Exposure measures the distribution of exposures of low percentile\\
    潜在未来敞口衡量的是低百分位数的敞口分布
    \item Expected exposure measures the mean distribution of exposures at a given future date prior to the maturity of the longest maturity exposure in the netting group\\
    期望敞口衡量的是在给定未来日期的敞口均值分布，时间范围在净额组合中最长期限之前
    \item The weights of expected positive exposure could represent the proportion that an individual expected exposures of the entire time interval
    期望正敞口的权重可以代表单个期望敞口在整个时间区间中的比例
\end{choices}

\begin{correctanswer}
    B
\end{correctanswer}

\begin{explanation}
    \textbf{A选项说法恰当。}
    \begin{itemize}
        \item 当前敞口是在交易对手违约时投资者可能损失的市场价值，假设破产时没有回收，取零和市场价值的较大值，因此不能选为"不准确"的选项。
    \end{itemize}

    \textbf{B选项说法不恰当。}
    \begin{itemize}
        \item 潜在未来敞口（PFE）衡量的是高百分位数的敞口分布（如95\%或99\%分位数），用于衡量极端情况下的敞口，而不是低百分位数的敞口分布。
    \end{itemize}

    \textbf{C选项说法恰当。}
    \begin{itemize}
        \item 期望敞口（Expected Exposure）衡量的是在给定未来日期的敞口均值分布，时间范围在净额组合中最长期限之前，因此不能选为"不准确"的选项。
    \end{itemize}

    \textbf{D选项说法恰当。}
    \begin{itemize}
        \item 期望正敞口（EPE）的权重可以代表单个期望敞口在整个时间区间中的比例，这是期望正敞口计算的基本原理，因此不能选为"不准确"的选项。
    \end{itemize}
\end{explanation}

\checklist{潜在未来敞口（PFE）的定义}



\begin{question}
    Which of the following statements best reflects the reason why a financial institution does not need to consider aggregating stresses to the expected positive exposure (EPE) with its loan portfolio?
    \\ 哪一项陈述最好反映了为什么金融机构不需要将其贷款组合的期望正敞口（EPE）与压力测试相结合？
\end{question}

\begin{choices}
    \item Loans are not sensitive to market variables\\
    贷款对市场变量不敏感
    \item Stresses to EPE are not sensitive to market variables\\
    EPE的压力测试对市场变量不敏感
    \item The EPE and the loan portfolio are negatively correlated\\
    EPE和贷款组合是负相关的
    \item The EPE and the loan portfolio are positively correlated\\
    EPE和贷款组合是正相关的
\end{choices}

\begin{correctanswer}
    A
\end{correctanswer}

\begin{explanation}
    \textbf{A选项正确。}
    \begin{itemize}
        \item 贷款对市场变量不敏感，因此不会因市场变量变化而产生风险敞口变化。EPE压力测试主要针对市场变量（如利率、汇率、股价等）的变化，而贷款组合的价值主要取决于借款人的信用质量，而不是市场变量，因此不需要考虑将EPE的压力测试与贷款组合结合。
    \end{itemize}

    \textbf{B选项错误。}
    \begin{itemize}
        \item EPE的压力测试对市场变量是敏感的，这是EPE压力测试的核心特征，因此这个选项与实际情况相反。
    \end{itemize}

    \textbf{C选项错误。}
    \begin{itemize}
        \item 相关性不是主要考虑因素。虽然EPE和贷款组合可能存在相关性，但关键原因是贷款对市场变量的不敏感性，而不是相关性本身。
    \end{itemize}

    \textbf{D选项错误。}
    \begin{itemize}
        \item 相关性不是主要考虑因素。无论EPE和贷款组合是正相关还是负相关，关键原因是贷款对市场变量的不敏感性，因此不需要考虑将EPE压力测试与贷款组合结合。
    \end{itemize}
\end{explanation}

\checklist{预期正敞口（EPE）压力测试与贷款组合风险管理的区别}



\begin{question}
    Is the following statement on the treatment of counterparty credit risk (CCR) correct?
    "Treating CCR as a market risk does not allow an institution to hedge market risk losses, and it exposes the institution to declines in counterparty creditworthiness and default."\\
    将交易对手信用风险（CCR）视为市场风险不允许机构对冲市场风险损失，并使机构暴露于交易对手信用质量下降和违约风险。
\end{question}

\begin{choices}
    \item The statement is correct with regard to both hedging market risk losses and counterparty creditworthiness and default\\
    关于对冲市场风险损失和交易对手信用质量下降和违约的说法是正确的
    \item The statement is incorrect with regard to both hedging market risk losses and counterparty creditworthiness and default\\
    关于对冲市场风险损失和交易对手信用质量下降和违约的说法是不正确的
    \item The statement is correct with regard to hedging market risk losses only\\
    关于对冲市场风险损失的说法是正确的
    \item The statement is correct with regard to counterparty creditworthiness and default only
    关于交易对手信用质量下降和违约的说法是正确的
\end{choices}

\begin{correctanswer}
    D
\end{correctanswer}

\begin{explanation}
    \textbf{A选项错误。}
    \begin{itemize}
        \item 题目陈述中"不允许机构对冲市场风险损失"是错误的。将CCR视为市场风险实际上允许对冲市场风险损失。
    \end{itemize}

    \textbf{B选项错误。}
    \begin{itemize}
        \item 题目陈述中"使机构暴露于交易对手信用质量下降和违约风险"是正确的，因此B选项说"陈述是错误的"不正确。
    \end{itemize}

    \textbf{C选项错误。}
    \begin{itemize}
        \item 题目陈述中"不允许机构对冲市场风险损失"是错误的，因此C选项说"陈述是正确的"不正确。
    \end{itemize}

    \textbf{D选项正确。}
    \begin{itemize}
        \item 题目陈述中"使机构暴露于交易对手信用质量下降和违约风险"是正确的。将CCR视为市场风险时，机构确实会暴露于交易对手信用质量下降和违约风险，因为信用风险无法通过市场风险对冲来完全消除。
    \end{itemize}
\end{explanation}

\checklist{压力测试中对手方违约概率与敞口的关系}


\begin{question}
    Risk managers at a commercial bank are preparing a stress test of the bank's capital ratios that assumes a severe financial market shock over a 2-year time horizon. The managers want to ensure that the stress test appropriately captures the impact of the stress scenario on the bank's investment portfolios, as well as on the bank's balance sheet items over the stress test horizon. Which of the following would be the most appropriate assumption or methodology for the managers to include in the stress test?
    \\ 一家商业银行的风险经理正在准备一项资本比率的压力测试，假设在未来2年的时间内发生严重的金融市场冲击。经理们希望确保压力测试适当地捕捉到压力情景对银行投资组合和资产负债表项目的影响。经理们应该在压力测试中包含以下哪一项假设或方法？
\end{question}

\begin{choices}
    \item Yields on speculative-grade corporate debt and short-term AAA-rated sovereign debt increase simultaneously\\
    投机级公司债和短期AAA级主权债收益率同时上升
    \item The default probabilities of derivatives counterparties increase as the bank's exposures to those counterparties increase\\
    衍生品交易对手的违约概率随着银行对这些交易对手的敞口增加而上升
    \item The volume of the bank's used auto loans falls much more than the volume of its new auto loans\\
    银行使用的二手车贷款量比新车贷款量下降更多
    \item The stress test is calibrated so that the bank is assumed to operate at the minimum capital ratio required to remain solvent\\
    压力测试是校准的，假设银行在维持偿付能力所需的最低资本比率下运营
\end{choices}

\begin{correctanswer}
    B
\end{correctanswer}

\begin{explanation}
    \textbf{A选项错误。}
    \begin{itemize}
        \item 投机级公司债和AAA级主权债收益率不会同时上升。在压力情景下，投资者会转向高质量资产（flight to quality），导致投机级债券收益率上升而高质量主权债收益率下降，这个选项不符合市场行为特征。
    \end{itemize}

    \textbf{B选项正确。}
    \begin{itemize}
        \item 衍生品交易对手的违约概率会随着银行对交易对手的敞口增加而上升，这反映了错路风险（wrong-way risk）。在压力情景下，敞口增加和对手方信用质量下降同时发生，导致整体风险增加，这是压力测试中需要重点考虑的风险特征。
    \end{itemize}

    \textbf{C选项错误。}
    \begin{itemize}
        \item 二手车贷款和新车贷款的下降幅度差异不是关键因素。压力测试应该关注系统性风险因素，而不是过于关注具体业务细节，这个选项不符合压力测试的整体性要求。
    \end{itemize}

    \textbf{D选项错误。}
    \begin{itemize}
        \item 压力测试不应该假设银行仅维持最低资本要求。这会导致测试结果过于乐观，不符合审慎监管要求。压力测试应该考虑更保守的资本水平。
    \end{itemize}
\end{explanation}

\checklist{压力测试中对手方违约概率与敞口的关系}







\end{document}